\newif\ifieeeaccessclass
\IfFileExists{ieeeaccess.cls}{%
\ieeeaccessclasstrue
\documentclass{ieeeaccess}
}{%
\ieeeaccessclassfalse
\documentclass[journal]{IEEEtran}
}
\usepackage{cite}
\usepackage{amsmath,amssymb,amsfonts}
\usepackage{algorithmic}
\usepackage{graphicx}
\usepackage{textcomp}
\usepackage{xurl}
\usepackage[hidelinks]{hyperref}
\DeclareUnicodeCharacter{2264}{$\leq$}
\DeclareUnicodeCharacter{2265}{$\geq$}
\setlength{\emergencystretch}{2em}
\hbadness=10000
\vbadness=10000
\ifieeeaccessclass\else
\providecommand{\history}[1]{}
\providecommand{\doi}[1]{}
\providecommand{\authorrefmark}[1]{\textsuperscript{#1}}
\newcommand{\address}[2][]{ }
\providecommand{\tfootnote}[1]{}
\providecommand{\corresp}[1]{}
\providecommand{\headeretal}{et al.}
\newdimen\titlepgskip
\newenvironment{keywords}{\begin{IEEEkeywords}}{\end{IEEEkeywords}}
\fi
\def\BibTeX{{\rm B\kern-.05em{\sc i\kern-.025em b}\kern-.08em
    T\kern-.1667em\lower.7ex\hbox{E}\kern-.125emX}}

\begin{document}

\history{Date of publication XXXX 00, 0000, date of current version XXXX 00, 0000.}
\doi{10.1109/ACCESS.XXXX.DOI}
\title{Cybersecurity and Cyber-Resilience for the European Cloud--Edge--IoT Computing Continuum: A Strategic Capability Roadmap}
\author{\uppercase{Author Name}\authorrefmark{1}, \uppercase{Author Name}\authorrefmark{2}}
\address[1]{Affiliation 1, City, Country (e-mail: author1@example.com)}
\address[2]{Affiliation 2, City, Country (e-mail: author2@example.com)}
\tfootnote{This manuscript is a LaTeX conversion of the source document into the IEEE Access journal template. Replace the placeholder metadata, affiliations, and funding statement as needed.}
\markboth{Author \headeretal: Cybersecurity and Cyber-Resilience for the European Cloud--Edge--IoT Computing Continuum}{Author \headeretal: Cybersecurity and Cyber-Resilience for the European Cloud--Edge--IoT Computing Continuum}
\corresp{Corresponding author: Author Name (e-mail: corresponding.author@example.com).}

\begin{abstract}
The growing integration of artificial intelligence across cloud, edge,
and IoT infrastructures makes cybersecurity, resilience, and digital
sovereignty an urgent strategic challenge for Europe. As computation and
intelligence become distributed across a heterogeneous computing
continuum, traditional perimeter-centric security models are no longer
sufficient. This challenge is particularly acute given the simultaneous
operationalisation of NIS2, the Cyber Resilience Act, the AI Act, and
DORA, together with increasing concern over Europe's structural
dependencies in cloud ecosystems, hardware supply chains, and critical
digital infrastructure. This paper addresses that gap by proposing a
structured methodology for identifying, profiling, and prioritising
cybersecurity capabilities for the European Cloud--Edge--IoT continuum.
The methodology proceeds through three analytical stages: architectural
decomposition into security layers and cross-cutting concerns,
capability formulation through a seven-dimension profiling template, and
priority clustering with phased transition matrices. Applying this
methodology, the paper identifies ten strategic capabilities spanning
hardware roots of trust, post-quantum cryptographic agility, lightweight
edge security, federated interoperability, supply chain verification,
cognitive security operations, pan-European network federation,
confidential computing, resilient infrastructure, and proactive threat
management. These capabilities are consolidated into three executive
priority clusters --- Sovereign Trusted Infrastructure, Cognitive
Security and Autonomous Cyber Defence, and Federated Operational
Sovereignty --- and analysed across a 2026--2040 transition horizon. The
paper further advances the concept of autonomous, federated, and
certified Security-by-Design as a strategic target for Europe, and
derives implementation-oriented recommendations linking research,
policy, industrial capability, and sovereignty objectives. It is
relevant both as a methodological contribution to cybersecurity
roadmapping and as a timely strategic framework for Europe's AI-enabled,
resilient, and sovereign digital future.
\end{abstract}

\begin{keywords}
Cloud--Edge--IoT continuum, cyber-resilience, cybersecurity roadmap, digital sovereignty, European Union, federated security, Security-by-Design.
\end{keywords}

\titlepgskip=-15pt
\maketitle

\section{Introduction}\label{introduction}

The ongoing evolution from centralised cloud infrastructures toward a
distributed Cloud--Edge--IoT continuum is fundamentally reshaping the
design, deployment, and operation of digital systems {[}42, 43, 44,
45{]}. By moving computation, storage, and intelligence closer to end
devices, users, and physical processes, this paradigm enables
low-latency services, real-time decision-making, operational autonomy,
and scalable data-driven innovation {[}46, 47, 48, 49, 50, 51{]} across
critical domains such as industry, mobility, energy, healthcare,
telecommunications, and public services. However, this architectural
shift also introduces a substantially more complex cybersecurity
landscape.

Unlike centralised environments, the Cloud--Edge--IoT continuum is
characterised by decentralisation, heterogeneity, physical exposure, and
the coexistence of multiple administrative and trust domains {[}64,
65{]}. Security can no longer rely primarily on static perimeters,
centralised visibility, or uniform infrastructure assumptions. Instead,
defenders must secure systems that are geographically dispersed,
operationally diverse, resource-constrained, and often exposed to both
cyber and physical threats {[}53, 66{]}. As a result, traditional
reactive and siloed security models are increasingly insufficient for
addressing the scale, dynamism, and interdependence of the continuum
{[}57, 67{]}.

This challenge is particularly important in the European context. Europe
has established a strong regulatory and normative foundation for
trustworthy digital systems through instruments such as NIS2 {[}1{]},
the Cyber Resilience Act {[}2{]}, the AI Act {[}3{]}, and the evolving
EU cybersecurity certification framework {[}60, 61{]}. NIS2 establishes
a unified legal framework across critical sectors and strengthens risk
management, reporting, cooperation, and enforcement, though
transposition has proven significantly delayed, with approximately 20 of
27 Member States having completed national implementation by early 2026
against a legal deadline of October 2024, and the Commission having
issued reasoned opinions to 19 Member States in May 2025 {[}4{]}. The AI
Act introduces a risk-based framework for trustworthy AI, with high-risk
system obligations under Articles 9, 13, and 15 entering applicability
from August 2026 under the default timetable, though the Digital Omnibus
on AI (COM(2025) 836) proposes deferral to December 2027 for certain
categories {[}5{]}. EUCC provides the first operational European
cybersecurity certification mechanism for ICT products {[}6{]}, while
EUCS for cloud services remains unadopted as of April 2026 --- a
structural policy impasse now likely to be resolved through the 2026
Cybersecurity Act revision (COM(2026) 13) rather than through the
original ENISA candidate scheme process {[}7{]}.

At the same time, Europe's position remains asymmetric. While it is
comparatively strong in trust, privacy, certification, and governance,
it still faces structural dependencies in cloud ecosystems, digital
infrastructure, and semiconductor supply chains. The European Chips Act
{[}8{]} explicitly links semiconductor capacity to technological
sovereignty and reduced external dependency, with the EU targeting 20\%
of global semiconductor production by 2030 --- though the European
Court of Auditors' Special Report 12/2025 projects a more realistic
trajectory of approximately 11.7\%, and the cancellation of Intel's
Magdeburg flagship fabrication facility in July 2025 has underscored the
fragility of this ambition {[}9{]}. Recent analyses continue to
highlight dependence on non-EU software and cloud providers, with
Synergy Research Group reporting that AWS, Azure, and GCP collectively
hold approximately 70\% of the European cloud market as of mid-2025
{[}10{]}, alongside broader concerns about strategic reliance on foreign
actors and single service providers.

Against this background, the strategic goal of this paper is to support
the transition from \emph{reactive, siloed security} toward
\emph{autonomous, federated, and certified Security-by-Design {[}56, 57,
58, 59{]}} for the Cloud--Edge--IoT continuum. In this paper,
\emph{autonomous} refers to infrastructures capable of sustaining
trustworthy and resilient operation under dynamic and adverse
conditions; \emph{federated} refers to interoperable security and trust
mechanisms spanning multiple providers, sectors, and jurisdictions,
like SSO systems; and \emph{certified Security-by-Design} refers to
embedding assurance, compliance, and trust requirements throughout the
lifecycle of systems, components, and services rather than applying them
retrospectively.

The main argument of the paper is that Europe's comparative advantage in
trust, privacy, governance, and certification can become strategically
meaningful only if it is translated into deployable cybersecurity
capabilities and coherent technical priorities. Therefore, this paper
approaches cybersecurity not merely as a narrow technical protection
problem, but as a combined challenge of architecture, operational
resilience, governance, interoperability, and sovereignty {[}60, 61, 62,
63{]}. On this basis, the paper identifies and structures the
cybersecurity priorities that are most critical for building a secure,
resilient, and sovereign European Cloud--Edge--IoT ecosystem.

The paper makes the following contributions. It frames
cybersecurity in the Cloud--Edge--IoT continuum as a multi-layer
challenge spanning technical security, operational resilience,
governance, interoperability, and digital sovereignty {[}52, 53, 54,
55{]}. It advances a strategic transition from reactive, siloed
protection models toward autonomous, federated, and certified
Security-by-Design {[}56, 57, 58, 59{]}. It proposes a three-stage
methodology --- architectural decomposition, seven-dimension
capability profiling, and priority clustering with transition matrices
--- that connects technical cybersecurity capabilities with broader
European strategic concerns including sovereignty, trust, and regulatory
alignment. It identifies and profiles ten strategic capabilities and
consolidates them into three priority clusters analysed across a
2026--2040 horizon. And it provides a prioritisation-oriented foundation
that can support research and innovation planning, policy alignment, and
deployment strategy in Europe.

The remainder of this paper is organised as follows. Section 2 presents
the background, state of the art, and key driving factors, with
particular attention to the European position. Section 3 introduces the
methodology. Section 4 presents the architectural layers and
corresponding focus areas. Section 5 analyses the ten identified
capabilities through the seven-dimension profiling template. Section 6
synthesises these capabilities into priority clusters and transition
matrices. Section 7 presents strategic recommendations. Section 8
discusses validation, limitations, and methodological positioning.
Section 9 concludes.

\section{Background, State of the Art, and Driving
Factors}\label{background-state-of-the-art-and-driving-factors}

The transition from centralised cloud computing toward a distributed
Cloud--Edge--IoT continuum is reshaping the cybersecurity landscape.
While this paradigm enables low-latency services, localised
intelligence, real-time analytics, and increased operational autonomy,
it also introduces a fundamentally more complex security environment. In
contrast to the relatively controlled conditions of centralised
hyperscale cloud infrastructures, the continuum is characterised by
large numbers of heterogeneous, geographically dispersed, and
resource-constrained nodes operating across multiple administrative and
trust domains {[}64, 65{]}. As a result, the attack surface expands
significantly, creating infrastructures that are, in practice, too
numerous, too diverse, and too dispersed to be secured through
traditional perimeter-oriented approaches alone {[}57, 67{]}.

A central challenge stems from decentralisation itself. Security models
developed for centralised cloud environments assume relatively stable
boundaries, consolidated monitoring, and strong physical and logical
control over infrastructure. These assumptions do not hold in edge and
IoT deployments, where services, workloads, and data increasingly move
closer to end users, devices, and operational environments. This
decentralisation complicates visibility, orchestration, patching,
identity management, and incident response, while also increasing the
likelihood of inconsistent security postures across the continuum.

A second major challenge is the physical exposure of edge
infrastructure. Unlike hyperscaler data centres, edge nodes are often
deployed in the field --- in factories, smart-city installations,
transport systems, telecom sites, or remote industrial environments.
These deployments rarely benefit from the same degree of physical
protection, controlled access, or environmental safeguards as
centralised facilities. Consequently, they are more vulnerable to device
tampering, hardware manipulation, fault injection, side-channel attacks,
and supply-chain compromise {[}53, 66{]}. This makes hardware-rooted
trust, secure boot, attestation, and tamper-aware system design
increasingly important as baseline requirements rather than optional
enhancements.

Simultaneously, the cybersecurity landscape is being shaped by an
unprecedented regulatory push. In Europe, the operationalisation of NIS2
{[}1{]}, the Cyber Resilience Act {[}2{]}, the AI Act {[}3{]}, and the
Digital Operational Resilience Act (DORA) {[}11{]} is creating immediate
pressure for organisations to demonstrate security, resilience,
accountability, and compliance across distributed and interconnected
infrastructures. The challenge is not only legal alignment, but
technical translation: regulatory obligations must be converted into
deployable security architectures, assurance mechanisms, lifecycle
controls, and governance models that can function across a fragmented
ecosystem of cloud, edge, and IoT stakeholders. The simultaneity of this
regulatory activation --- with CRA enforcement beginning September 2026
and full application in December 2027, AI Act high-risk provisions from
August 2026, DORA fully applied from January 2025, and the Data Act
{[}12{]} operational from September 2025 --- creates a compliance
coordination challenge that disproportionately burdens SMEs and smaller
Member States.

These technical and regulatory pressures are closely linked to the
broader strategic objective of digital sovereignty. As Europe expands
its digital infrastructure and data-driven services, systemic dependence
on non-EU providers for key hardware, platforms, and security-critical
components introduces both geopolitical and operational risk.
Extra-territorial data access concerns --- underscored by Microsoft
France's June 2025 confirmation before the French Senate that it cannot
guarantee EU data will never be accessed under the US CLOUD Act {[}13{]}
--- foreign technology dependencies, and vendor lock-in can undermine
Europe's capacity to guarantee autonomy, trust, and resilience in
critical digital infrastructures.

\subsection{Europe in Relation to the Current State of the
Art}\label{europe-in-relation-to-the-current-state-of-the-art}

In relation to the global state of the art, Europe holds a distinctive
position. Its comparative advantage lies less in hyperscale platform
dominance and more in its leadership in trust, privacy, governance, and
certification {[}60, 61, 62{]}. Europe has established a world-leading
regulatory and normative foundation through instruments such as the GDPR
{[}14{]}, NIS2, the CRA, and the AI Act. Together, these frameworks
embed privacy-by-design, trust-by-design, and accountability into the
legal and operational fabric of digital systems. In parallel, European
certification initiatives --- notably EUCC, which became operational on
27 February 2025 under Commission Implementing Regulation (EU) 2024/482
{[}6{]} --- provide a rigorous basis for assurance and trusted
evaluation. The European Cybersecurity Competence Centre (ECCC)
{[}15{]}, the national coordination centres network, and the Horizon
Europe funding ecosystem create institutional support for long-term
research, innovation, and deployment. Initiatives such as the Simpl
middleware programme {[}16{]} and the Common European Data Spaces
further demonstrate Europe's ambition to build interoperable, federated,
and trustworthy digital infrastructures.

However, these strengths coexist with significant structural gaps.
Europe remains critically dependent on non-EU secure silicon, including
trusted platform modules, secure enclaves, and other hardware roots of
trust. While EU-headquartered vendors --- Infineon, NXP,
STMicroelectronics --- design the dominant TPM products and have
achieved EUCC `High' certification for multiple secure element families
(STMicroelectronics ST54J/K certified by ANSSI on 1 April 2025; NXP
SN300/JCOP in July--August 2025; Infineon IFX\_CCI in December 2025
{[}6{]}), fabrication depends on non-EU foundries (primarily TSMC) and
the designs remain proprietary, preventing sovereign auditability at the
silicon level. Intel, AMD, and ARM control the entire operational
trusted execution environment market. Threat intelligence and
operational cyber defence capabilities remain fragmented across national
and sectoral boundaries, with no unified pan-European cyber situational
awareness capability at scale. The enterprise SIEM market is dominated
by US platforms --- Splunk/Cisco, Microsoft Sentinel, and Palo Alto
Networks (which acquired IBM QRadar SaaS in August 2024) --- with no
EU-headquartered vendor in the Gartner Leaders quadrant as of October
2025 {[}17{]}. Europe also lacks production-scale PQC-capable hardware
tailored for edge and IoT environments, though the NIS Cooperation
Group's ``Coordinated Implementation Roadmap for the Transition to
Post-Quantum Cryptography'' (v1.1, 23 June 2025) {[}18{]} now provides a
binding Member State migration framework with milestones at 2026, 2030,
and 2035.

Taken together, this position suggests that Europe is not starting from
weakness, but from asymmetry. It is strong in governance, trust,
assurance, and regulatory design, yet weaker in industrial scale,
unified operationalisation, and sovereign control of critical enabling
technologies. This gap between normative leadership and deployable
capability --- what we term the \emph{ambition--implementation gap} ---
defines the central challenge for the next phase of European
cybersecurity strategy in the Cloud--Edge--IoT continuum.

\section{System model of the Cloud-Edge-IoT
ecosystem}\label{system-model-of-the-cloud-edge-iot-ecosystem}

These are

Vertical Stack:

\begin{itemize}
\item
  \textbf{Layer 1: Hardware-Centric Security}
\item
  \textbf{Layer 2: Software-Centric Security (Middleware/OS)}
\item
  \textbf{Layer 3: Application-Centric Security (Apps/Services)}
\item
  \textbf{Layer 4: User \& Orchestration-Centric Security}
\item
  \textbf{Layer 5: Data-Centric Security}
\item
  \textbf{Layer 6: Network Infrastructure \& Perimeter-Centric Security}
\end{itemize}

The Transversal Verticals (Cross-Layers):

\begin{itemize}
\item
  \textbf{Cross-Layer 1: Automated Security Processes \& Cyber-Resilient
  AI Lifecycle}
\item
  \textbf{Cross-Layer 2: Dynamic Risk Management \& Operational Cyber
  Resilience}
\item
  \textbf{Cross-Layer 3: Interoperable Security Mechanisms \&
  Open-Source Verifiability}
\item
  \textbf{Cross-Layer 4: Compliance-by-Design, Agile Governance \&
  Automated Assurance}
\end{itemize}

\section{Methodology}\label{methodology}

The methodology presented in this section is designed to transform a
broad and heterogeneous strategic landscape --- spanning hardware
security, post-quantum cryptography, AI-driven defence, federated
governance, and supply chain integrity across the European
cloud--edge--IoT computing continuum --- into a structured roadmap for
research, innovation, policy, deployment, and ecosystem development. The
analytical process proceeds through three staged phases ---
architectural decomposition, capability formulation and profiling, and
priority clustering with phased progression --- followed by an iterative
validation step.

Two foundational commitments shape the methodology. First, it is
intended to support strategic roadmap construction, not deterministic
forecasting. In complex socio-technical ecosystems, progress is
non-linear and depends on interacting technical, regulatory,
operational, and market developments whose timing and interdependencies
cannot be predicted with precision. Second, the methodology is designed
to connect technical cybersecurity capabilities with broader European
strategic concerns --- resilience, interoperability, trust, and what we
term \emph{Sovereignty Level}: the degree to which critical
capabilities, control points, and dependencies remain under credible
European influence, governance, or operational control. Sovereignty
Level is employed throughout the methodology as an interpretive
strategic lens rather than as a rigid scoring model; it guides
assessment of whether progress in cybersecurity and cyber-resilience
strengthens Europe's strategic control, trusted autonomy, and reduced
dependency in critical parts of the computing continuum.

The methodological process is structured into four interconnected steps:

\begin{itemize}
\item
  \textbf{Stage 1. Core Dimensions and Focus Areas Identification}:
  which establishes the analytical structure
\item
  \textbf{Stage 2. Capability Identification and Profiling:} which
  translates problem domains into strategic capabilities evaluated
  through a consistent analytical template
\item
  \textbf{Stage 3. Priority Clustering and Transition Matrix:} which
  groups capabilities into broader strategic priorities and positions
  them across phased pathways
\item
  \textbf{Validation step:} drawing on expert review, consultation,
  and iterative refinement.
\end{itemize}

\subsection{Stage 1: Core Dimensions and Focus Area
Identification}\label{stage-1-core-dimensions-and-focus-area-identification}


\begin{figure*}[t]
\centering
\IfFileExists{media/image1.png}{%
\includegraphics[width=\textwidth]{media/image1.png}
}{%
\fbox{\parbox{0.95\textwidth}{\centering Placeholder: media/image1.png not found.}}
}
\caption{Figure 1: Security Framework with separated security verticals in dedicated boxes}
\label{fig:security-framework}
\end{figure*}


The first stage provides the analytical scope and structure of the
roadmap by decomposing the cloud--edge--IoT computing continuum into a
set of architectural vertical and cross-cutting layers. These dimensions
are treated analytically distinct but operationally interdependent: a
vulnerability or governance gap at one layer typically has consequences
that propagate to adjacent layers, and effective security requires
coordinated action across the full vertical and horizontal extent of the
architecture.

The purpose of this stage is to answer three questions. Where are the
major security and resilience problem domains within the distributed
computing continuum? At what architectural layers or cross-layers do
these problems emerge? And which focus areas are sufficiently strategic,
cross-cutting, or underdeveloped to justify dedicated roadmap attention?

Focus areas are derived from four complementary sources. The first is
the architectural characteristics of distributed cloud--edge--IoT
systems themselves: the heterogeneity of execution environments spanning
HPC, cloud, edge, and device tiers creates structurally distinct
security requirements at each level of the stack. The second is the set
of recurring cybersecurity and cyber-resilience challenges observed
across these tiers --- hardware integrity, software supply chain
verification, workload isolation, identity federation, data
confidentiality, and network perimeter control. The third comprises
cross-cutting strategic concerns that cannot be localised to a single
layer: sovereignty, trust, interoperability, resilience, assurance, and
regulatory alignment. The fourth is the practical need to connect
technical architecture with operational and governance realities ---
recognising that a capability with no credible deployment pathway or
regulatory anchor is unlikely to achieve ecosystem-wide impact
regardless of its technical merit.

The output of Stage 1 is not a priority list but a structured map of the
problem space, organised into six Security Layers (the vertical stack)
and four Transversal Cross-Layers (the horizontal enablers), as detailed
in Section 4. Within this structured map, some focus areas acquire
greater strategic importance because they influence Europe's future
Sovereignty Level --- particularly where dependencies on external
providers, fabrication capacity, or assurance mechanisms are
structurally significant. Stage 1 identifies these sovereignty-sensitive
focus areas explicitly so that subsequent stages can track their
evolution as a strategic concern alongside purely technical maturity.

\subsection{Stage 2: Capability Identification and
Profiling}\label{stage-2-capability-identification-and-profiling}

Once the architectural map and its constituent focus areas are defined,
the methodology moves from problem framing to capability formulation. A
capability, in the context of this roadmap, is understood as a strategic
ability that the European ecosystem should be able to design, deploy,
govern, standardise, assure, coordinate, or scale in order to achieve
trustworthy, resilient, and sovereign operation of the computing
continuum. Capabilities are not limited to technical controls; they may
also encompass governance mechanisms, interoperability instruments,
assurance functions, ecosystem coordination capacities, and deployment
enablers.

Capabilities are derived by translating each focus area identified in
Stage 1 into actionable strategic needs, guided by a consistent
question: what must Europe be able to do --- technically,
institutionally, and industrially --- to address this focus area at the
scale and with the assurance level that the regulatory and threat
environment demands? The ten capabilities identified through this
process (detailed in Section 5) span the full range from hardware roots
of trust and post-quantum cryptographic agility to federated security
interoperability and proactive cyber threat management.

\subsubsection{The Seven-Dimension Capability
Profile}\label{the-seven-dimension-capability-profile}

Each capability is assessed through a seven-dimension analytical
template applied identically across all priorities. This template
ensures that every capability is characterised not only from a technical
perspective, but also in terms of its regulatory alignment, investment
pathway, failure modes, and contribution to European strategic autonomy.
The consistency of the template enables structured cross-priority
comparison --- a feature essential both for funding prioritisation and
for identifying shared dependencies and bottlenecks that span multiple
capabilities.

\textbf{Dimension 1: Baseline and Target Objective.} This dimension
grounds the capability in reality by describing the present condition of
the European ecosystem, the intended future condition, and the strategic
change the capability is meant to achieve. The baseline captures current
industrial capacity, the degree of reliance on non-European vendors,
technical and economic barriers to adoption, and the state of
standardisation and certification. The target defines measurable
end-states --- quantified market share objectives, certification levels,
migration coverage metrics, or deployment thresholds.

\textbf{Dimension 2: Primary Drivers.} This dimension articulates the
external pressures that make the capability urgent and non-deferrable.
It distinguishes between threat drivers (adversarial capabilities or
attack vectors that are already active or imminently scalable) and
regulatory drivers (specific EU legal obligations with enforcement
timelines, citing the precise directive or regulation and its relevant
provisions).

\textbf{Dimension 3: High-Impact Activities.} This dimension identifies
three to five concrete, actionable research and innovation activities
expected to yield the greatest strategic return on investment. Each
activity is tagged with its primary target phase (Phase 1, Phase 2, or
Phase 3).

\textbf{Dimension 4: Catalysts and Enablers.} This dimension identifies
the concrete mechanisms, funding vehicles, and regulatory
standardisations required to accelerate the capability from baseline
toward target. Three categories are distinguished: R\&I and funding
alignment (Horizon Europe clusters, Digital Europe Programme, ECCC
Strategic Agenda), regulatory and standardisation (CEN/CENELEC, ETSI,
ISO standards), and ecosystem enablers (pilot infrastructures,
industrial partnerships, SME capacity-building).

\textbf{Dimension 5: Disruptors and Systemic Risks.} This dimension
identifies risks that could render planned deployments obsolete or
unachievable: technical cliffs (defensive capabilities rendered
ineffective by adversarial breakthroughs), supply chain and market risks
(fabrication bottlenecks, vendor resistance, monopolisation), and
geopolitical risks (export controls, extra-territorial jurisdiction
conflicts).

\textbf{Dimension 6: Failsafe Mechanisms and Resilience.} This dimension
specifies engineering fallbacks if a capability's primary mechanism is
compromised, delayed, or only partially realised: hybrid transitional
schemes, graceful degradation modes, autonomous isolation protocols, and
implementation diversity mandates.

\textbf{Dimension 7: Transversal Meta-Layer (Policy Alignment).} This
dimension connects the capability to the EU's overarching strategic
autonomy agenda across three axes: sovereignty and autonomy, trust and
ethics, and sustainable security.

\subsection{Stage 3: Priority Clustering and Transition
Matrix}\label{stage-3-priority-clustering-and-transition-matrix}

Not every capability can or should be treated as a standalone roadmap
priority. Capabilities that share architectural proximity, common
enabling conditions, similar policy needs, overlapping ecosystem actors,
or convergent strategic outcomes benefit from coordinated treatment.
Stage 3 therefore groups the ten capabilities into three broader
Priority Clusters:

\begin{itemize}
\item
  \textbf{Priority Cluster 1: Sovereign Trusted Infrastructure} --- the
  hardware, cryptographic, and resilience foundations that anchor trust
  across the continuum.
\item
  \textbf{Priority Cluster 2: Cognitive Security and Autonomous Cyber
  Defence} --- the AI-driven operational layer that scales defence
  beyond human-paced response.
\item
  \textbf{Priority Cluster 3: Federated Operational Sovereignty} --- the
  interoperability, supply-chain, and federation substrate that enables
  a unified European trust fabric.
\end{itemize}

Each cluster is analysed through a three-phase transition matrix:

\begin{itemize}
\item
  \textbf{Phase 1: Near-Term Implementation and Pilot Actions
  (2026--2030)} --- deploying near-ready solutions, enforcing enacted
  regulation, and closing governance gaps before technical work scales.
\item
  \textbf{Phase 2: Mid-Term Scale-Up, Industrial R\&D and
  Standardisation (2030--2035)} --- industrialising sovereign
  capabilities, federating architectures, and building the workforce and
  institutional infrastructure the technical work requires.
\item
  \textbf{Phase 3: Long-Term Strategic Autonomy, Frontier R\&I and
  Leadership (2035--2040)} --- achieving full operational maturity,
  pioneering frontier research, and resolving structural doctrine and
  oversight questions.
\end{itemize}

\subsubsection{Analytical
Perspectives}\label{analytical-perspectives}

For each priority cluster and phase, the assessment considers four
analytical perspectives: \emph{R\&I maturity} (technical readiness and
industrialisation feasibility), \emph{policy, regulation, and standards}
(legislative enactment, enforcement credibility, standards availability,
conformity assessment capacity), \emph{operational deployment and
adoption} (movement beyond laboratory environments into operational
use), and \emph{market uptake and strategic impact} (commercial
traction, European provider competitiveness, measurable Sovereignty
Level improvement).

\subsubsection{Transition Conditions}\label{transition-conditions}

This roadmap uses \emph{transition conditions} rather than hard gates
between phases. Transition conditions are evidence-informed progression
markers indicating when a priority cluster appears sufficiently prepared
to advance. They are framed as analytical signals of increasing
readiness, coherence, and ecosystem support --- not binary thresholds or
deterministic requirements. This design choice reflects the non-linear
nature of progress in complex socio-technical ecosystems, where
adversarial breakthroughs, regulatory delays, fabrication bottlenecks,
and geopolitical disruptions can alter timelines unpredictably.

A transition condition is a short statement indicating that a technical
basis has become sufficiently mature for the next stage of deployment; a
policy or standardisation basis has emerged that reduces adoption
uncertainty; a deployment pathway has been validated through pilots; a
strategic ecosystem condition has materialised; or a Sovereignty Level
improvement has become plausible through reduced dependency or increased
European capacity.

Transition conditions are not promises, deterministic predictions, or
precise quantitative thresholds unless evidence genuinely supports them.
They should be interpreted as structured indicators of directional
readiness --- useful for strategic planning and coordination, but
subject to reassessment as the ecosystem evolves.

\section{Core Dimensions and Focus
Areas}\label{core-dimensions-and-focus-areas}

The roadmap is structured across a cohesive architectural framework
spanning four continuum environments --- HPC, Cloud, Edge, and Devices.
The capabilities are organised into six Security Layers (the vertical
stack) and four Transversal Cross-Layers (the horizontal enablers). This
section describes each dimension and its constituent focus areas.

\subsection{Security Layers (Vertical
Stack)}\label{security-layers-vertical-stack}

\textbf{Layer 1: Hardware-Centric Security.} This layer addresses the
physical and silicon-level trust foundations of the computing continuum.
It encompasses two focus areas. \emph{Layer 1.1 (Hardware Roots of
Trust)} concerns cryptographically verifiable supply chain trust
ensuring physical device integrity from fabrication to deployment, with
a strategic focus on driving European open-hardware designs (RISC-V)
toward EUCC `High' assurance certification {[}6{]}. \emph{Layer 1.2
(Physical-Layer Threat Detection and Fingerprinting)} exploits intrinsic
wireless channel characteristics --- Channel State Information, Carrier
Frequency Offset, and RSSI --- to establish real-time device
authentication and anomaly detection at the lowest communication layer,
complementing upper-layer cryptographic protections.

\textbf{Layer 2: Software-Centric Security (Middleware/OS).} This layer
secures the software substrate on which continuum workloads execute.
\emph{Layer 2.1 (OS and Middleware Hardening)} addresses the security of
hypervisors, container runtimes, and distributed operating systems, with
emphasis on sovereign, lightweight Linux distributions optimised for
edge resource constraints and reduced attack surfaces. \emph{Layer 2.2
(Secure OTA Patch and Update Management)} recognises that edge nodes
deployed in the field --- smart grids, O-RAN radios, industrial gateways
--- require cryptographically verified, failure-resilient over-the-air
update pipelines to maintain security posture without physical access.

\textbf{Layer 3: Application-Centric Security (Apps/Services).} This
layer protects the workloads and service interfaces that execute across
the continuum. \emph{Layer 3.1 (Workload Security)} secures
containerised applications and serverless functions through hardened
service meshes and sidecar proxy architectures providing mutual TLS and
granular traffic inspection. \emph{Layer 3.2 (In-Band Network Telemetry
and Deep Observability)} embeds smart telemetry directly into data
packets traversing containerised workloads, enabling real-time security
monitoring without dedicated out-of-band infrastructure. \emph{Layer 3.3
(API Security Gateway and Policy Enforcement)} provides a unified
enforcement point for authentication tokens, rate limiting, schema
validation, and payload inspection as containerised AI workloads
increasingly communicate through REST/gRPC APIs.

\textbf{Layer 4: User and Orchestration-Centric Security.} This layer
governs identity, access, and policy enforcement across the
multi-provider continuum. \emph{Layer 4.1 (Zero-Trust Orchestration)}
implements Policy Enforcement Point / Policy Decision Point
architectures enabling continuous authorisation and dynamic access
controls adapted to request context, aligned with NIST SP 800-207
{[}19{]}. \emph{Layer 4.2 (Federated Identity and Access Management)}
utilises Self-Sovereign Identity principles for cryptographically
secure, decentralised authentication of machines, users, and automated
workloads. \emph{Layer 4.3 (Privileged Access Management and
Just-In-Time Access)} addresses the lateral movement vector behind major
APT campaigns by replacing over-provisioned standing privileges with
time-bounded, contextually justified access grants.

\textbf{Layer 5: Data-Centric Security.} This layer provides end-to-end
data protection across the continuum lifecycle. \emph{Layer 5.1
(Confidential Computing)} delivers data-in-use security through Trusted
Execution Environments spanning the Cloud-to-IoT data path. \emph{Layer
5.2 (Post-Quantum Cryptography)} addresses the migration to
quantum-resistant algorithms --- ML-KEM (FIPS 203), ML-DSA (FIPS 204),
and SLH-DSA (FIPS 205), finalised by NIST in August 2024
{[}20{]}{[}21{]}{[}22{]} --- for resource-constrained devices to
neutralise ``Harvest Now, Decrypt Later'' threats. \emph{Layer 5.3
(Signalling Data Security)} embeds trust directly into network
signalling architecture. \emph{Layer 5.4 (Distributed Reputation and
Trust Networking)} implements real-time reputation systems built on
distributed ledger technology. \emph{Layer 5.5 (Quantum Key Distribution
for Critical Links)} provides information-theoretically secure key
exchange for high-value point-to-point links. \emph{Layer 5.6 (Secure
Multi-Party Computation)} enables multiple mutually distrusting parties
to jointly compute functions over combined private inputs without
revealing those inputs.

\textbf{Layer 6: Network Infrastructure and Perimeter-Centric Security.}
This layer secures the communication fabric of the continuum.
\emph{Layer 6.1 (Telco-cloud and O-RAN Security)} hardens distributed
5G/6G cores and RAN components in alignment with the SNS JU SRIA
{[}23{]}, ensuring sovereign control over network slice isolation.
\emph{Layer 6.2 (Physical Resilience and Safety)} adapts IT/OT gateway
security for high-availability edge environments to prevent digital
breaches from causing kinetic damage. \emph{Layer 6.3 (Edge-Native
Autonomous Guardians)} hosts lightweight, localised security control
functions directly at far-edge gateways or within Open-RAN Radio
Intelligent Controllers. \emph{Layer 6.4 (Non-Terrestrial Network and
Satellite Security)} defines security requirements for 6G NTN
integration aligned with 3GPP NTN specifications {[}24{]} and the EU
Space Programme.

\subsection{Transversal Cross-Layers (Horizontal
Enablers)}\label{transversal-cross-layers-horizontal-enablers}

\textbf{Cross-Layer 1: Automated Security Processes and Cyber-Resilient
AI Lifecycle.} This cross-layer spans the full lifecycle of software and
AI component security. \emph{CL1.1 (Continuous DevSecOps)} embeds
AI-driven code analysis and automated vulnerability remediation at the
earliest stages of the software lifecycle. \emph{CL1.2 (Model Integrity
and DataBOMs)} secures AI training pipelines against data poisoning and
unauthorised model manipulation. \emph{CL1.3 (Dynamic Supply Chain
Verification and Runtime SBOMs)} implements dynamic Software Bills of
Materials and cryptographic proof of execution to continuously verify
the integrity of containerised microservices and AI components as they
migrate across the multi-provider continuum. \emph{CL1.4 (Automated
Adversarial Testing and Continuous Purple Teaming)} provides the
adversarial counterpart to DevSecOps: continuous automated penetration
testing, fuzzing, and purple-team automation against live staging
environments.

\textbf{Cross-Layer 2: Dynamic Risk Management and Operational Cyber
Resilience.} This cross-layer governs real-time threat detection,
analytics, and operational response. \emph{CL2.1 (Automated SOC and
Threat Intelligence)} enables SOC federation across Member States using
legally protected platforms to overcome institutional trust barriers
{[}25{]}. \emph{CL2.2 (High-Precision Telemetry and Security
Signalling)} implements smart telemetry to feed real-time threat
intelligence without overwhelming network bandwidth. \emph{CL2.3
(Internet of Behaviour and Evidence Sharing)} extracts and shares
behavioural patterns across the infrastructure. \emph{CL2.4 (AI-Native
Analytics Integration via NWDAF and LLMSecOps)} routes cybersecurity
telemetry through centralised analytics hubs using 5G NWDAF architecture
{[}26{]} and large language model-based security operations. \emph{CL2.5
(Cyber Threat Intelligence Lifecycle Management)} governs the full CTI
lifecycle as a prerequisite for the federated SOC vision. \emph{CL2.6
(Deception Technology and Active Defence)} deploys decoy assets ---
honeypots, honeynets, credential canaries, fake API endpoints --- that
generate high-fidelity, low-noise alerts upon adversary interaction.

\textbf{Cross-Layer 3: Interoperable Security Mechanisms and Open-Source
Verifiability.} This cross-layer enables secure interoperation across
heterogeneous providers and jurisdictions. \emph{CL3.1 (Interoperability
and Simpl Integration)} mandates integration into the Simpl open-source
middleware {[}16{]} to unlock Common European Data Spaces. \emph{CL3.2
(Neutral Host and Multi-Tenant Isolation)} enforces Zero-Trust edge
operations across different network partitions under infrastructure
sharing models. \emph{CL3.3 (Pan-European NaaS Federation and
Cross-Border Edge Roaming)} establishes standardised interfaces
leveraging CAMARA Network APIs {[}27{]} for a unified
Network-as-a-Service platform. \emph{CL3.4 (Semantic Security
Interoperability)} ensures participating organisations share common
ontologies for security events, asset taxonomies, and risk
classifications. \emph{CL3.5 (Cyber Range and Digital Twin Federation)}
federates national cyber ranges into shared EU-wide validation
infrastructure aligned with the Cyber Solidarity Act.

\textbf{Cross-Layer 4: Compliance-by-Design, Agile Governance and
Automated Assurance.} This cross-layer embeds regulatory compliance into
technical architecture. \emph{CL4.1 (Post-Quantum Readiness and
Zero-Trust Architectures)} embeds continuous, AI-driven Zero-Trust
verification natively into cloud-edge infrastructure. \emph{CL4.2
(Automated Certification)} replaces point-in-time audits with real-time
certification tracking for cloud services (EUCS, when adopted) and edge
products (EUCC) {[}6{]} throughout their lifecycle. \emph{CL4.3
(Multicloud Key Management as a Service)} manages cryptographic keys
across dynamic multi-provider computing continuum operations.
\emph{CL4.4 (Privacy-Preserving Federated Computation)} enables
cross-jurisdictional AI training and inference over distributed private
datasets. \emph{CL4.6 (Cyber Risk Quantification and Board-Level
Reporting)} translates technical security posture into financial and
operational risk language for decision-makers. \emph{CL4.7 (AI
Trustworthiness and Algorithmic Accountability)} addresses the
trustworthiness of AI systems that themselves make autonomous security
decisions, in compliance with AI Act Article 9 requirements {[}3{]}.

\section{Capability Profiles}\label{capability-profiles}

The ten strategic capabilities identified in this roadmap are analysed
through the seven-dimension profiling template described in Section
3.2.1. This section presents each capability profile in full.

\subsection{P1: Hardware Roots of Trust and EUCC `High'
Certification}\label{p1-hardware-roots-of-trust-and-eucc-high-certification}

\textbf{Baseline and Target.} The current European landscape reflects a
structurally dependent but industrially active hardware security
ecosystem. EU-headquartered vendors --- Infineon, NXP,
STMicroelectronics --- design the dominant secure element products
deployed globally, and all three achieved EUCC `High' certification for
multiple product families in 2025: STMicroelectronics ST54J/K (ANSSI, 1
April 2025), NXP SN300/JCOP (July--August 2025), and Infineon IFX\_CCI
G12 (December 2025) {[}6{]}. However, fabrication depends on non-EU
foundries (primarily TSMC), and designs remain proprietary, preventing
sovereign auditability at the silicon level. No EU-origin RISC-V
processor has received EUCC or Common Criteria certification at any
assurance level as of April 2026. Intel, AMD, and ARM control the entire
operational TEE market. The target is ubiquitous deployment of
EU-designed, open-source secure elements on EU-controlled fabrication
capacity, certified at EUCC `High', providing a verifiable trust
foundation for the entire continuum. By 2040, at least 60\% of critical
infrastructure edge nodes should utilise EU-designed, EUCC `High'
certified secure elements.

\textbf{Primary Drivers.} \emph{Threat:} state-sponsored exploitation of
JTAG debug interfaces and side-channel attacks on field-deployed edge
nodes; supply chain insertion of hardware backdoors. \emph{Regulatory:}
CRA Articles 10--11 imposing hardware security and lifecycle management
requirements with enforcement from December 2027 {[}2{]}; NIS2 Article
21(2)(e) mandating cryptographic policies for essential entities
{[}1{]}; AI Act Articles 9 and 15 requiring hardware integrity
guarantees for high-risk AI systems {[}3{]}.

\textbf{High-Impact Activities.} (Phase 1) Fund 3--5 RISC-V secure
element pilots under Horizon Europe Cluster 3 targeting EUCC `High'
pre-evaluation; submit OpenTitan {[}28{]} and Keystone {[}29{]}
prototypes for certification assessment. (Phase 2) Bring the first EUCC
`High' certified EU-manufactured RISC-V secure elements into production
via EU Chips Act pilot lines, including NanoIC (imec, inaugurated
February 2026) {[}8{]} and FAMES (CEA-Leti). (Phase 3) Establish
EU-designed RISC-V secure elements as the hardware default across all
new EU ICT products.

\textbf{Catalysts and Enablers.} EU Chips Act pilot fabrication lines
{[}8{]}; ECCC Strategic Agenda {[}15{]} funding for open-hardware
security; CEN/CENELEC JTC 13 hardware security standards; the DARE SGA1
RISC-V initiative (EuroHPC, launched March 2025, €240M) {[}30{]};
Giesecke+Devrient's OpenTitan CC certification path.

\textbf{Disruptors and Systemic Risks.} EU Chips Act fabrication
capacity shortfall --- ECA Special Report 12/2025 projects only
\textasciitilde11.7\% global semiconductor share by 2030, against the
20\% target {[}9{]}. Intel Magdeburg cancellation (July 2025) removed
the largest EU-based advanced node commitment. TSMC's Dresden ESMC joint
venture (28/22nm, production expected late 2027) does not address the
leading-edge nodes required for advanced secure elements. Side-channel
attack breakthroughs rendering current enclave architectures obsolete.

\textbf{Failsafe Mechanisms.} Design all pilot architectures with
explicit TEE abstraction layers to avoid lock-in to Intel TDX, AMD
SEV-SNP, or ARM TrustZone during the Phase 1 dependency period. Deploy
software-based isolation (TEEs) as a secondary layer if specific
hardware enclaves are compromised. Implement autonomous node quarantine
via eBPF-based network policy enforcement if attestation fails.

\textbf{Transversal Meta-Layer.} \emph{Sovereignty:} eliminates
proprietary hardware dependencies and enables sovereign silicon
auditability --- the current dependency level is effectively 100\% for
TEE infrastructure and \textasciitilde100\% for fabrication.
\emph{Trust:} EUCC `High' provides the highest assurance level under the
EU cybersecurity certification framework, aligned with Common Criteria
EAL4+ and AVA\_VAN.5. \emph{Sustainability:} energy-efficient secure
silicon designs for far-edge deployment reduce the carbon footprint of
cryptographic operations.

\subsection{P2: Post-Quantum Cryptographic Agility for
Edge/IoT}\label{p2-post-quantum-cryptographic-agility-for-edgeiot}

\textbf{Baseline and Target.} Current edge and IoT environments depend
almost entirely on classical cryptographic algorithms (RSA, ECDH, ECDSA)
vulnerable to future quantum computers. Cryptographic asset inventories
are largely absent, and migration readiness is weak outside government
and financial sectors. The NIS Cooperation Group's PQC roadmap (v1.1,
June 2025) {[}18{]} establishes a binding EU migration framework:
national roadmaps and cryptographic inventories by December 2026;
high-risk critical infrastructure transitioned by December 2030; full
completion by December 2035. NIST FIPS 203 (ML-KEM), FIPS 204 (ML-DSA),
and FIPS 205 (SLH-DSA) were finalised in August 2024
{[}20{]}{[}21{]}{[}22{]}; FIPS 206 (FN-DSA/Falcon) remains
pre-publication as of April 2026; HQC was selected as a backup KEM in
March 2025 {[}31{]}. The target is a quantum-native continuum with full
cryptographic agility enabling algorithm updates without hardware
replacement.

\textbf{Primary Drivers.} \emph{Threat:} the Q-Day horizon (current
expert consensus: 10--20 years {[}32{]}, with significant lower-bound
uncertainty) and active ``Harvest Now, Decrypt Later'' campaigns.
\emph{Regulatory:} NIS CG PQC roadmap {[}18{]} with 2030/2035
milestones; ANSSI requirement for PQC capability in all Qualification
submissions from 2027; BSI TR-02102-1 v2026-01 approving ML-KEM, ML-DSA,
SLH-DSA, HQC, FrodoKEM, and Classic McEliece.

\textbf{High-Impact Activities.} (Phase 1) Map cryptographic inventories
per NIS CG roadmap requirements; optimise FIPS 203/204/205 algorithms
for 8/16-bit microcontrollers; deploy hybrid ML-KEM + ECDH-P384 in
cloud-side implementations per ETSI TS 103 744 v1.2.1 {[}33{]}. (Phase
2) Mandate hybrid PQC in all new Industrial IoT and financial sector
deployments requiring DORA compliance; validate lightweight PQC for
sub-100 µJ zero-energy nodes under 6G ISAC testbeds. (Phase 3) Establish
PQC-native silicon as the EU procurement default.

\textbf{Catalysts and Enablers.} NIST FIPS 203/204/205
{[}20{]}{[}21{]}{[}22{]}; ETSI TS 103 744 v1.2.1 hybrid key exchange
profiles {[}33{]}; ENISA PQC migration guidelines; ECCC Strategic Agenda
{[}15{]}; the IETF LAMPS working group X.509 certificate drafts (in RFC
Editor Queue); open-source crypto-agility frameworks (e.g., liboqs,
Bouncy Castle PQC).

\textbf{Disruptors and Systemic Risks.} Early CRQC arrival before
migration is complete. Performance overhead of PQC on battery-powered
IoT devices (ML-KEM demonstrates \textasciitilde1.5× computational
overhead versus ECDH at equivalent security levels on ARM Cortex-M4
{[}32{]}). Insufficient cryptographic inventory of legacy OT/IoT systems
preventing migration planning. Interoperability gaps between hybrid
implementations.

\textbf{Failsafe Mechanisms.} Deploy hybrid cryptographic schemes
(classical + PQC) during the transition phase to maintain backward
compatibility while providing harvest resistance. Implement
crypto-agility frameworks enabling algorithm switching within sub-second
latency if vulnerabilities are discovered. Diversify across PQC
algorithm families (lattice-based, code-based, hash-based) to avoid
single-point cryptanalytic failures.

\textbf{Transversal Meta-Layer.} \emph{Sovereignty:} protects long-term
European governmental and industrial secrets from future quantum
decryption --- the 18 Member State joint statement ``Securing Tomorrow,
Today'' (November 2024) formalises this as a collective sovereignty
concern. \emph{Trust:} alignment with GDPR Article 32 requirement for
state-of-the-art security measures. \emph{Sustainability:} lightweight
PQC implementations reduce the energy overhead and extend device
lifecycles, avoiding premature hardware replacement.

\subsection{P3: Lightweight Cryptography and Physical-Layer Security
for the Extreme
Edge}\label{p3-lightweight-cryptography-and-physical-layer-security-for-the-extreme-edge}

\textbf{Baseline and Target.} Extreme edge environments --- massive
machine-type communication nodes, body-area networks, zero-energy
backscatter radios --- remain constrained by limited compute, energy,
memory, and connectivity, making conventional cryptographic protections
impractical. The target is a class of lightweight, adaptive security
mechanisms combining efficient cryptography with physical-layer
primitives tailored for highly constrained environments, enabling
quantum-resistant authentication for devices operating at sub-milliwatt
power budgets.

\textbf{Primary Drivers.} \emph{Threat:} the proliferation of mMTC
requiring 100+ devices per m³ with severe energy constraints targeting
40-year battery lifespans or energy-harvesting operation.
\emph{Regulatory:} CRA applicability to products with digital elements,
including resource-constrained IoT devices.

\textbf{High-Impact Activities.} (Phase 1) Test group-based
authentication and early lightweight PQC algorithms on
resource-constrained IoT; validate RF fingerprinting beyond controlled
lab environments. (Phase 2) Hardware-software co-design for
neuromorphic/PIM AI-assisted security on EU open-hardware platforms.
(Phase 3) Deploy zero-energy physical-layer authentication as a standard
identity mechanism for battery-less mMTC nodes.

\textbf{Catalysts and Enablers.} NIST lightweight cryptography efforts;
6G and O-RAN research programmes under SNS JU {[}23{]}; ETSI SmartBAN
standardisation for body-area networks; European experimentation
platforms for edge and wireless security.

\textbf{Disruptors and Systemic Risks.} Variable reliability of
PHY-based techniques under environmental changes; device heterogeneity;
danger of deploying under-validated mechanisms in safety-critical
systems; immaturity of deployment standards.

\textbf{Failsafe Mechanisms.} Layered combinations of lightweight
cryptography and higher-layer verification; adaptive confidence
thresholds; dynamic fallback to localised trust boundaries using purely
physical-layer signatures when upper-layer cryptographic processing
exceeds the available energy budget.

\textbf{Transversal Meta-Layer.} \emph{Sovereignty:} secures the far
edge of European digital infrastructure where conventional approaches
cannot reach. \emph{Sustainability:} energy-aware security is a
precondition for scalable extreme-edge deployment without
disproportionate environmental cost.

\subsection{P4: Security Interoperability via Simpl
Middleware}\label{p4-security-interoperability-via-simpl-middleware}

\textbf{Baseline and Target.} European cloud-edge ecosystems remain
fragmented across proprietary vendor silos with incompatible security
policies, authentication mechanisms, and data exchange protocols. The
Simpl programme --- funded through a €41M development contract to the
Sovereign-X consortium (led by Eviden Belgium, awarded February 2024)
and a €65M second tender to InfrateX (led by Sopra Steria + NTT DATA)
--- has produced an iterative MVP available since January 2025 with
enhanced releases in June and December 2025, hosted on EU GitLab
(code.europa.eu/simpl/simpl-open) {[}16{]}. The target is a unified
European ``Trust Fabric'' enabling secure, federated interoperability
across heterogeneous providers through standardised security Minimal
Interoperability Mechanisms (MIMs) within the Simpl stack. However,
OASC's MIM6 (Security management) remains at ``Work item'' status with
no formally published specification as of April 2026.

\textbf{Primary Drivers.} \emph{Regulatory:} Data Act (Regulation (EU)
2023/2854, applicable from 12 September 2025) {[}12{]} mandating
interoperability and portability; NIS2 requirements for coordinated
incident response across federated infrastructure. \emph{Market:}
breaking vendor lock-in --- the primary barrier to EU provider
competitiveness is proprietary API ecosystems, not technical
superiority.

\textbf{High-Impact Activities.} (Phase 1) Launch Simpl security MIM
pilots in at least three Common European Data Spaces (Manufacturing,
Health, Energy). (Phase 2) Deploy automated policy translation via
standardised Simpl MIMs; achieve adoption as ratified CEN/CENELEC
standards referenced in the EU Cloud Rulebook. (Phase 3) Establish
Simpl-based security as the universal standard for Data Space operation.

\textbf{Catalysts and Enablers.} Digital Europe Programme funding for
Simpl development; Gaia-X Trust Framework 3.0 ``Danube'' (November
2025); CEN-CLC/JTC 25 (established September 2024, chaired by Sebastian
Steinbuss/IDSA) for data-space standards; Eclipse Dataspace Protocol
nearing v1.0 release; IPCEI-CIS ICRA Reference Architecture v2.0
(January 2026).

\textbf{Disruptors and Systemic Risks.} Vendor resistance to open
standards; weakest-link security failures in federated systems where one
provider's compromise cascades; insufficient incentives for
interoperability; uneven maturity across sectors and Member States.

\textbf{Failsafe Mechanisms.} Capability to autonomously unfederate or
isolate nodes that fail to meet the shared security baseline; Zero-Trust
architectures where federation does not imply implicit trust;
cryptographic attestation of policy enforcement.

\textbf{Transversal Meta-Layer.} \emph{Sovereignty:} breaking
hyperscaler lock-in (AWS, Azure, GCP hold \textasciitilde70\% of the EU
cloud market {[}10{]}; European-headquartered providers collectively
hold \textasciitilde15\%, stable since 2022 but down from 29\% in 2017).
\emph{Trust:} transparent, verifiable security mechanisms enabling
cross-border trust.

\subsection{P5: Supply Chain Verification and Secured
Infrastructure}\label{p5-supply-chain-verification-and-secured-infrastructure}

\textbf{Baseline and Target.} Current digital infrastructures depend on
complex, opaque, and globally distributed supply chains. CRA Annex I
Part II will mandate static SBOM disclosure from December 2027 {[}2{]};
NIS2 Article 21(2)(h) makes supply chain risk assessment mandatory for
essential entities {[}1{]}; DORA Article 28 requires continuous ICT
third-party risk monitoring {[}11{]} --- operationalised through the
first 19 Critical ICT Third-Party Provider designations on 18 November
2025 {[}34{]}. The target is verifiable and continuously assessable
supply chains spanning hardware, software, firmware, and service
dependencies, progressing from static disclosure (Phase 1) through
runtime verification (Phase 2) to real-time ZKP-attested provenance
(Phase 3).

\textbf{Primary Drivers.} \emph{Threat:} rising software supply chain
attacks (SolarWinds-class); disaggregation of telco clouds introducing
extreme multi-vendor complexity. \emph{Regulatory:} simultaneous CRA
SBOM, NIS2 supply chain risk assessment, and DORA ICT third-party
monitoring obligations.

\textbf{High-Impact Activities.} (Phase 1) Submit CRA-mandated static
SBOM declarations; initiate Digital Product Passport pilot programmes;
advance ZKP verification toward demonstrator readiness. (Phase 2) Deploy
Runtime SBOMs with cryptographic proof-of-execution in critical
infrastructure; validate ZKP batch verification operationally. (Phase 3)
Operationalise ZKP-verified supply chain attestation at real-time
transaction scale.

\textbf{Catalysts and Enablers.} CRA standardisation request M/606
{[}2{]}; CEN-CLC/JTC 13/WG 9 horizontal standards (EN 40000 series); EU
Ecodesign Regulation Digital Product Passport requirements for ICT
hardware; DORA Delegated Regulation (EU) 2024/1502 for CTPP assessment.

\textbf{Disruptors and Systemic Risks.} Global concentration of
manufacturing dependencies; limited multi-tier auditability; high
complexity of continuously verifying dynamic supply chains; SME
exclusion from compliance pathways.

\textbf{Failsafe Mechanisms.} Graceful isolation of unverified
components into quarantined network slices; multi-vendor
diversification; trusted fallback components; continuous monitoring for
supply chain anomalies.

\textbf{Transversal Meta-Layer.} \emph{Sovereignty:} verifiability tied
to accountability and procurement governance. \emph{Trust:} the DORA
CTPP designation of 19 providers (including AWS, Google Cloud,
Microsoft, SAP, Deutsche Telekom, Orange, SWIFT {[}34{]}) demonstrates
the operational reach of EU supply chain oversight.

\subsection{P6: Cognitive SecOps (LLMSecOps) and Automated Threat
Intelligence}\label{p6-cognitive-secops-llmsecops-and-automated-threat-intelligence}

\textbf{Baseline and Target.} Security operations remain heavily manual,
fragmented, and overloaded by alert volumes and analyst scarcity.
US-platform SIEMs dominate the EU market, with no EU-headquartered
vendor in the Gartner Magic Quadrant Leaders quadrant (October 2025)
{[}17{]}. Foundation models for SOC copilot functions are US-origin. The
most concrete European sovereign security AI initiative is Thales'
GenAI4AID, built on Mistral AI models; the OpenEuroLLM consortium
(launched February 2025, €37.4M budget) has no dedicated security model
stream {[}35{]}. The target is a cognitive SecOps model where
EU-sovereign language models and AI-driven reasoning systems augment
detection, triage, investigation, and response while remaining
controllable, auditable, and secure under AI Act Article 9 conformity
assessment.

\textbf{Primary Drivers.} \emph{Threat:} weaponisation of generative AI
for polymorphic malware, automated exploit discovery, and social
engineering at scale. \emph{Regulatory:} AI Act Articles 9, 13, and 15
creating mandatory requirements for autonomous security AI classified as
high-risk systems {[}3{]}; a 24--36 month regulatory gap where
obligations precede conformity assessment capacity.

\textbf{High-Impact Activities.} (Phase 1) Integrate commercial LLMs
into SOC analyst workflows for alert triage and MITRE ATT\&CK technique
mapping; initiate training of at least one EU-sovereign SLM (≤10B
parameters) on GDPR-compliant security corpora using EuroHPC compute.
(Phase 2) Deploy EU-sovereign SLMs at the edge for real-time SOC
analytics without cloud inference dependency; publish a
CEN/CENELEC-standardised adversarial robustness test suite. (Phase 3)
Deploy swarm-based agentic AI for decentralised threat intelligence.

\textbf{Catalysts and Enablers.} EuroHPC AI Factories --- 19 confirmed
across 16 Member States with €2.7B cumulative funding; IT4LIA (CINECA),
PIAST (Poznań), and MeluXina-AI (Luxembourg) explicitly list
cybersecurity as a vertical {[}36{]}. ETSI GS ZSM 009 (closed-loop
automation) {[}37{]}; ECCC Strategic Agenda; Sekoia.io/HarfangLab French
Open XDR Platform consortium.

\textbf{Disruptors and Systemic Risks.} Hallucinations and automation
bias; adversarial prompt manipulation; data leakage into models; no
standardised adversarial robustness test suite for security AI (the
single most consequential standards gap in this cluster); overreliance
on opaque AI behaviour in critical contexts.

\textbf{Failsafe Mechanisms.} Human-in-the-loop validation for Tier-1
critical actions; bounded autonomy with ETSI GS ZSM 009 governance;
``shadow-mode'' rollback if autonomous agents exhibit unstable control
loops; audit logging and policy-constrained orchestration.

\textbf{Transversal Meta-Layer.} \emph{Sovereignty:} reducing
\textasciitilde80\% US-platform dependency in enterprise SIEM.
\emph{Trust:} AI Act compliance-by-design, explainability, and
accountability for autonomous security decisions.

\subsection{P7: Pan-European NaaS Federation (``Schengen for
Cloud-Edge/AI'')}\label{p7-pan-european-naas-federation-schengen-for-cloud-edgeai}

\textbf{Baseline and Target.} European digital infrastructures remain
fragmented across providers, jurisdictions, and operational models.
CAMARA's Fall 2025 meta-release comprises 60 APIs (10 stable) {[}27{]};
GSMA Open Gateway has reached 86 operator groups representing 300+
networks and 80\% of global mobile connections {[}38{]}. CEF Digital
Call 4 selected 12 5G corridor projects in Q4 2025. However,
cross-border edge compute roaming does not exist at operational scale.
The target is a trusted pan-European NaaS federation enabling secure,
policy-aligned workload migration across operator boundaries.

\textbf{Primary Drivers.} \emph{Threat:} the ``Sovereignty Cliff'' ---
lack of a true Digital Single Market prevents European operators from
scaling to compete with US/China hyperscalers. \emph{Regulatory:} the
Digital Networks Act proposal (COM(2026) 16, adopted January 2026)
signals EC intent for network federation, though it does not mandate
specific API standards.

\textbf{High-Impact Activities.} (Phase 1) Deploy harmonised CAMARA APIs
in 2--3 bilateral 5G corridor agreements. (Phase 2) Scale cross-border
roaming to ≥5 EU operator pairs; operationalise ``Data Visa'' mechanisms
for dynamic AI model migration. (Phase 3) Operate the EU NaaS federation
as a commercially viable unified jurisdiction for data workloads.

\textbf{Catalysts and Enablers.} IPCEI-CIS (€2.6B total, 19 direct
participants, 12 Member States) {[}39{]}; CAMARA consortium {[}27{]};
SNS JU (€630M cumulative EU funding, 100 projects) {[}23{]}; Aduna
(Ericsson/operator JV) as primary aggregator.

\textbf{Disruptors and Systemic Risks.} Inconsistent governance across
Member States; operator misalignment; hyperscaler platform dominance; US
CLOUD Act jurisdictional conflict.

\textbf{Failsafe Mechanisms.} Graceful degradation protocols redirecting
workloads to national infrastructure; regional autonomy zones;
multi-operator fallback routing.

\textbf{Transversal Meta-Layer.} \emph{Sovereignty:} the ``Schengen for
Cloud'' legal framework would resolve structural CLOUD Act conflict.
\emph{Trust:} legal interoperability and equitable distribution of
strategic infrastructure value.

\subsection{P8: Confidential Computing and Privacy-Preserving AI
Pipelines}\label{p8-confidential-computing-and-privacy-preserving-ai-pipelines}

\textbf{Baseline and Target.} Sensitive data processing across
cloud-edge-AI pipelines remains exposed to platform-level trust
assumptions. Intel TDX, AMD SEV-SNP, and ARM TrustZone dominate all
operational TEE deployments; EU-origin alternatives (RISC-V Keystone
{[}29{]}, OpenTitan {[}28{]}) remain at research/prototype stage. The
target is end-to-end privacy-preserving AI pipelines using TEEs,
homomorphic encryption, and verifiable processing guarantees,
progressively migrating to EU-origin hardware as Cluster 1 matures.

\textbf{Primary Drivers.} \emph{Regulatory:} GDPR Article 25 (data
protection by design) {[}14{]}; AI Act trustworthiness requirements;
EUCS (when adopted) certification for confidential cloud services.
\emph{Market:} protection of industrial IP during federated training
across multi-provider Data Spaces.

\textbf{High-Impact Activities.} (Phase 1) Design pilot architectures
with TEE abstraction layers; initiate FHE benchmarking for edge-viable
use cases. (Phase 2) Deploy RISC-V Keystone TEEs in multi-provider Data
Spaces alongside Intel TDX infrastructure. (Phase 3) Operationalise FHE
end-to-end from IoT sensor to HPC compute node.

\textbf{Catalysts and Enablers.} European Data Space programmes; EUCS
certification development; Chips JU RISC-V TEE development; GDPR-aligned
data governance frameworks.

\textbf{Disruptors and Systemic Risks.} Hardware dependence on
non-European TEE vendors throughout Phase 1; performance overhead of
FHE; side-channel vulnerabilities in current TEE implementations;
limited cross-vendor TEE attestation interoperability.

\textbf{Failsafe Mechanisms.} Hybrid privacy architectures; selective
workload confinement; limit processing to local federated learning nodes
if edge TEE integrity cannot be verified.

\textbf{Transversal Meta-Layer.} \emph{Sovereignty:} Intel/AMD/ARM TEE
dependency is \textasciitilde100\% for operational deployments --- this
is a critical sovereignty gap. \emph{Trust:} verifiable data protection
strengthens AI adoption legitimacy.

\subsection{P9: Resilient-by-Design Infrastructure and Autonomous
``Island
Mode''}\label{p9-resilient-by-design-infrastructure-and-autonomous-island-mode}

\textbf{Baseline and Target.} Many critical digital infrastructures
assume persistent connectivity, stable central coordination, and
continuous service dependencies. The target is resilient-by-design
infrastructure capable of autonomous local operation, graceful
degradation, and secure continuation under disruption through ``Island
Mode'' --- sustaining a minimum of 72 hours of autonomous operation
under network partition or orchestration-plane compromise.

\textbf{Primary Drivers.} \emph{Threat:} escalating cyberattacks,
extreme weather, and geopolitical tensions causing cascading failures.
\emph{Regulatory:} NIS2 business continuity mandates for essential
entities {[}1{]}; DORA Article 11 business continuity requirements for
financial entities {[}11{]}; SNS JU SRIA requirements for 6G resilience
and AI-RAN {[}23{]}.

\textbf{High-Impact Activities.} (Phase 1) Validate Island Mode
soft-state protocols in 5G testbeds aligned with SNS JU objectives;
develop autonomous fallback mechanisms for edge nodes lacking
survivability modes. (Phase 2) Achieve standardised certification of
72-hour Island Mode under 3GPP/ETSI fallback API profiles in EU 5G
cross-border corridor pilots. (Phase 3) Mandate validated 72-hour
autonomous Island Mode across all NIS2-essential and DORA-critical edge
deployments.

\textbf{Catalysts and Enablers.} DORA {[}11{]}; SNS JU 6G resilience
projects (HORSE, ROBUST-6G, ELASTIC) {[}23{]}; civil protection and
critical infrastructure programmes; resilience-oriented public
procurement.

\textbf{Disruptors and Systemic Risks.} Centralised architectural
legacies; poor failover design; excessive cloud dependence; insufficient
integration between safety, security, and operational resilience
engineering.

\textbf{Failsafe Mechanisms.} This capability inherently depends on
redundant trust anchors, segmented autonomy zones, offline continuity
functions, and tamper-resistant local decision logic. Zero-energy backup
layers or localised LPWAN micro-slices maintain a ``minimal viable
network'' for life-saving telemetry if standard edge power and computing
fail.

\textbf{Transversal Meta-Layer.} \emph{Sovereignty:} sovereign
infrastructure must be dependable under crisis. \emph{Resilience:} P9
directly operationalises the ``cyber-resilience'' concept at the
infrastructure level.

\subsection{P10: Proactive Cyber Threat Management and Defence
against
APTs}\label{p10-proactive-cyber-threat-management-and-defence-against-apts}

\textbf{Baseline and Target.} Threat management remains largely
reactive, with limited anticipation of coordinated, stealthy, and
persistent adversaries. Average APT dwell times remain at 150--200 days
in many European organisations. The target is a proactive defence
posture combining predictive threat intelligence, attack path analysis,
adversary emulation mapped to the MITRE ATT\&CK framework, and adaptive
Moving Target Defence, achieving dwell times below 72 hours in Phase 2
and below 24 hours in Phase 3 across all NIS2-essential deployments.

\textbf{Primary Drivers.} \emph{Threat:} nation-state actors blending
cyber espionage with organised crime; exploitation of complex supply
chains (SolarWinds-class attacks); stealthy techniques including
memory-only execution and abuse of legitimate administrative tools.
\emph{Regulatory:} NIS2 Article 21(2)(b) incident handling requirements
{[}1{]}.

\textbf{High-Impact Activities.} (Phase 1) Apply ML-based User and
Entity Behaviour Analytics for APT IoC detection; target dwell time
reduction below 72 hours in at least one certified NIS2-essential
entity. (Phase 2) Deploy deception infrastructure alongside AI detection
for model calibration; orchestrate autonomous triage under ETSI GS ZSM
009 governance {[}37{]}. (Phase 3) Operationalise Moving Target Defence
continuously shuffling IP addressing, network topology, and software
diversity.

\textbf{Catalysts and Enablers.} ENISA threat intelligence efforts;
CSIRT and SOC cooperation structures; MITRE ATT\&CK framework
integration into incident management; pan-European CTI sharing.

\textbf{Disruptors and Systemic Risks.} Intelligence fragmentation; poor
information-sharing incentives; adversarial adaptation outpacing
defensive innovation; organisational reluctance to invest in proactive
capabilities before incidents occur.

\textbf{Failsafe Mechanisms.} Autonomous real-time micro-segmentation
(Zero-Trust) to quarantine compromised components even if the specific
exploit is unknown; layered detection logic; deception and containment
zones; rapid recovery playbooks.

\textbf{Transversal Meta-Layer.} \emph{Sovereignty:} European strategic
preparedness requires sovereign threat intelligence independent of
US-origin feeds. \emph{Governance:} information sharing requires
trustworthy automation and proportionate handling of intelligence-driven
interventions.

\section{Priority Clusters and Transition
Matrix}\label{priority-clusters-and-transition-matrix}

The ten capabilities are consolidated into three executive priority
clusters analysed across two dimensions per phase: (i) Technical Actions
and Deployment Targets and (ii) Governance, Policy and Standards
Milestones. Each phase carries a distinct mandate reflecting the
characteristic transition from regulatory enforcement (Phase 1) through
industrial scale-up (Phase 2) to sovereign leadership (Phase 3).

\subsection{Priority Cluster 1: Sovereign Trusted
Infrastructure}\label{priority-cluster-1-sovereign-trusted-infrastructure}

This cluster establishes the cryptographically verifiable trust
foundation for the European cloud--edge--IoT continuum. It unifies five
capabilities: P1 (Hardware Roots of Trust), P2 (PQC Agility), P3
(Lightweight Cryptography and PhySec), P8 (Confidential Computing), and
P9 (Resilient-by-Design Infrastructure). This cluster provides the
hardware and cryptographic assurances that Clusters 2 and 3
progressively consume as their own capabilities mature.

The cluster addresses a single architectural premise: federated security
policies (Cluster 3) and autonomous defence capabilities (Cluster 2)
gain operational credibility in proportion to the maturity of the
execution substrate on which they run. That substrate must guarantee
three properties: hardware integrity verifiable through open-source
silicon roots of trust; quantum-resistant confidentiality through
crypto-agile schemes aligned with FIPS 203/204/205
{[}20{]}{[}21{]}{[}22{]}; and autonomous resilience enabling continued
operation under network or system-level compromise.

The shared critical dependency across this cluster is European sovereign
fabrication capacity under the EU Chips Act {[}8{]}. Hardware Roots of
Trust, PQC Agility, Lightweight Cryptography, and Confidential Computing
all require RISC-V-based secure elements manufactured on EU-controlled
process nodes to reach their Phase 2 targets. This fabrication
bottleneck --- underscored by the ECA's contested 20\% target {[}9{]}
and Intel Magdeburg's cancellation --- is the single most consequential
cross-priority risk in the roadmap.

\subsubsection{Phase 1: Foundations and Pilot Actions
(2026--2030)}\label{phase-1-foundations-and-pilot-actions-20262030}

\emph{Sovereignty Level: High Foreign Dependency}

\textbf{Technical Actions.} Fund RISC-V secure element projects
targeting EUCC `High' under Horizon Europe Cluster 3; target Fragmented
Piloting by 2030. Validate Island Mode soft-state protocols in 5G
testbeds aligned with SNS JU SRIA objectives {[}23{]}. Optimise FIPS
203/204/205 algorithms for 8/16-bit microcontrollers. Submit Keystone
{[}29{]} and OpenTitan {[}28{]} prototypes for EUCC pre-evaluation.
Validate RF fingerprinting in representative field conditions. Initiate
FHE benchmarking for edge-viable use cases, recognising that operational
deployment is not achievable before Phase 2. Establish 3--5 Horizon
Europe consortia operating RISC-V secure element pilots in controlled
critical infrastructure demonstrators. Design all pilot architectures
with explicit TEE abstraction layers to avoid lock-in to Intel TDX, AMD
SEV-SNP, or ARM TrustZone, which will dominate all operational TEE
deployments throughout this phase. Initiate PQC deployment in cloud-side
software implementations; far-edge IoT will run classical cryptography
exclusively throughout Phase 1. Develop autonomous fallback mechanisms
for edge nodes lacking survivability modes meeting DORA {[}11{]} or NIS2
{[}1{]} business continuity thresholds.

\textbf{Governance and Standards.} Align all R\&I outputs with the
regulatory enforcement timeline: NIS2 Article 21(2)(e) cryptographic
policies (legal obligation since October 2024, though transposition
remains incomplete in approximately one-third of Member States {[}4{]});
CRA Articles 10--11 hardware security certification (conformity
assessment body notification from June 2026, full application December
2027 {[}2{]}); AI Act Articles 9 and 15 hardware integrity requirements
(applicability from August 2026 under the default timetable, potentially
deferred to December 2027 under the Digital Omnibus {[}5{]}). Adopt the
NIS CG PQC roadmap {[}18{]} as the baseline for all new cryptographic
deployments. Engage in CEN/CENELEC JTC 13 hardware security standards
drafting. Prioritise achieving EUCC `High' certification for an
EU-origin RISC-V secure element before end of phase --- this is the
critical regulatory gate. Note that EUCC `High' has already been
achieved by Infineon, NXP, and STMicroelectronics for proprietary secure
elements {[}6{]}; the Phase 1 target is certification of an
\emph{open-source, EU-auditable} design. Initiate EU Chips Act
fabrication investment targeting pre-production milestones. Establish
strategic stockpiling and second-source agreements as interim mitigation
against foundry access disruption.

\subsubsection{Phase 2: Scale-Up, Industrial R\&D and Standardisation
(2030--2035)}\label{phase-2-scale-up-industrial-rd-and-standardisation-20302035}

\emph{Sovereignty Level: Emerging EU Ecosystem}

\textbf{Technical Actions.} Bring the first EUCC `High' certified
EU-manufactured RISC-V secure elements into production. Deploy in energy
substations and telecom critical infrastructure in at least five Member
States. Target Standardised Integration for Hardware RoT and Island Mode
by 2035. Validate hybrid ML-KEM + ECDH-P384 key encapsulation at
Industrial IoT scale. Demonstrate lightweight PQC for sub-100 µJ
zero-energy nodes under 6G ISAC testbeds. Deploy RISC-V Keystone TEEs
operationally in multi-provider Data Spaces alongside Intel TDX
infrastructure, validating cross-vendor TEE attestation
interoperability. Achieve standardised certification of 72-hour Island
Mode under 3GPP/ETSI fallback API profiles in at least three EU 5G
cross-border corridor pilots under SNS JU governance. Mandate hybrid PQC
schemes as standard in all new Industrial IoT and financial sector
deployments requiring DORA compliance.

\textbf{Governance and Standards.} Secure the definitive Phase 2 policy
gate: EUCC `High' certification of the first EU-manufactured RISC-V
secure element, transforming applied research into a market-accessible,
legally certifiable product. Adopt NIST FIPS 203/204/205 into ETSI TS
profiles for EU deployment contexts. Press ENISA to publish a mandatory
PQC migration timeline referenced in NIS2 sector-specific implementing
acts. Drive CEN/CENELEC hardware security standards through
ratification. Target 15--30\% market share for EU-designed secure
elements in new critical infrastructure procurement. Establish
second-source fabrication agreements between EU foundries.

\subsubsection{Phase 3: Strategic Autonomy and Sovereign Leadership
(2035--2040)}\label{phase-3-strategic-autonomy-and-sovereign-leadership-20352040}

\emph{Sovereignty Level: Full Sovereign Autonomy}

All five priorities target Ubiquitous Operation. Establish EU-designed
RISC-V secure elements as the hardware default across all new EU ICT
products, with 60\% or more of critical infrastructure edge nodes using
EUCC `High' certified EU-designed elements. Operationalise FHE
end-to-end from IoT sensor to HPC compute node. Deploy zero-energy
physical-layer authentication via RF fingerprinting as a standard
identity mechanism for battery-less mMTC nodes. Mandate validated
72-hour autonomous Island Mode survivability across all NIS2-essential
and DORA-critical edge deployments. Make EUCC `High' certification
mandatory for all new EU ICT products without waiver. Establish
PQC-native silicon as the EU procurement default. Achieve adoption of at
least one EU hardware security standard in ISO/IEC or ITU-T
international standardisation. R\&I investment orients toward
neuromorphic secure compute, post-second-generation quantum hardware,
and post-CMOS security architectures.

\subsection{Priority Cluster 2: Cognitive Security and Autonomous
Cyber
Defence}\label{priority-cluster-2-cognitive-security-and-autonomous-cyber-defence}

This cluster operationalises artificial intelligence as both the primary
instrument and a critical attack surface within the European cyber
defence architecture. It encompasses P6 (LLMSecOps) and P10 (Proactive
APT Management), together with cross-layer dimensions governing AI
system integrity: Model Integrity and DataBOMs (CL1.2), Dynamic Supply
Chain Verification (CL1.3), and AI Trustworthiness and Algorithmic
Accountability (CL4.7).

The cluster is structured around three inseparable sub-dimensions.
First, AI as a security instrument: deploying sovereign European
foundation models, SLMs, and edge-native agents for autonomous SOC
operations and real-time Moving Target Defence against APT campaigns
mapped to the MITRE ATT\&CK framework. Second, AI as an adversary
amplifier: recognising that generative capabilities are simultaneously
available for polymorphic malware generation and automated exploit
discovery. Third, AI as target: securing training pipelines and
inference endpoints against data poisoning, model evasion, and prompt
injection, in direct compliance with AI Act Article 9 requirements for
high-risk system robustness {[}3{]}.

As Cluster 1 matures its TEE and hardware attestation capabilities
through each phase, Cluster 2 progressively deploys autonomous agents
onto increasingly trustworthy substrates --- from cloud-hosted copilots
(Phase 1) to edge-native autonomous guardians (Phase 3). Delays in any
one cluster constrain the operational credibility of the other two at
each subsequent phase gate.

\subsubsection{Phase 1: Foundations and Pilot Actions
(2026--2030)}\label{phase-1-foundations-and-pilot-actions-20262030-1}

\emph{Sovereignty Level: High Reliance on US Foundation Models}

\textbf{Technical Actions.} Integrate commercial LLMs into SOC analyst
workflows for alert triage and MITRE ATT\&CK technique mapping; target
Siloed AI Copilots by 2030. Route NWDAF telemetry pipelines into
centralised AI analytics in 3--5 5G SA pilot deployments {[}26{]}. Apply
ML-based User and Entity Behaviour Analytics for APT IoC detection.
Target APT dwell time reduction below 72 hours in at least one certified
pilot deployment in a named NIS2-essential entity. Develop adversarial
robustness testing methodologies for prompt injection, model evasion,
and DataBOM verification; no operational certification framework exists
at phase start. Acknowledge that no certified autonomous containment
agent will reach operational deployment during Phase 1.

\textbf{Governance and Standards.} Prepare for AI Act Articles 9, 13,
and 15 (applicability from August 2026): mandatory requirements for
autonomous security AI classified as high-risk. Address the 24--36 month
regulatory gap where obligations precede conformity assessment body
operational readiness {[}5{]}. Adopt ETSI GS ZSM 009 {[}37{]} as the
governance reference for all autonomous SOC pilots. Operationalise NIS2
Article 21(2)(b) incident handling across pilot SOC federations.
Prioritise development of a standardised adversarial robustness test
suite under CEN/CENELEC or ETSI TC CYBER --- the single most
consequential standards gap in this cluster. Initiate training of at
least one EU-sovereign SLM (≤10B parameters) on GDPR-compliant security
corpora using EuroHPC compute {[}36{]}.

\subsubsection{Phase 2: Scale-Up, Industrial R\&D and Standardisation
(2030--2035)}\label{phase-2-scale-up-industrial-rd-and-standardisation-20302035-1}

\emph{Sovereignty Level: Sovereign EU Edge LLMs Emerging}

\textbf{Technical Actions.} Deploy EU-sovereign SLMs trained on MITRE
ATT\&CK-aligned corpora at the edge for real-time SOC analytics without
cloud inference dependency. Enable Guardian nodes to execute autonomous
Zero-Trust reconfigurations per NIST SP 800-207 {[}19{]}. Orchestrate
autonomous incident triage under ETSI GS ZSM 009 governance. Target APT
dwell time ≤72 hours in certified critical infrastructure deployments.
Deploy deception infrastructure (honeynets, credential canaries)
alongside AI detection for model calibration ground truth.

\textbf{Governance and Standards.} Secure the definitive Phase 2 policy
gate: operationalisation of a standardised AI Act conformity assessment
framework for autonomous security agents. Publish a
CEN/CENELEC-standardised adversarial robustness test suite. Formally
reference MITRE ATT\&CK in NIS2 sector-specific incident handling
guidance. Activate EuroHPC AI Factory compute contracts for EU-sovereign
security model training {[}36{]}. Target EU-sovereign SLMs covering
≥30\% of critical infrastructure SOC deployments in regulated sectors.

\subsubsection{Phase 3: Strategic Autonomy and Sovereign Leadership
(2035--2040)}\label{phase-3-strategic-autonomy-and-sovereign-leadership-20352040-1}

\emph{Sovereignty Level: Global Leadership in Trustworthy AI Security}

All sub-dimensions target Ubiquitous Operation. Deploy swarm-based
agentic AI defending the EU Cloud-Edge Continuum via self-evolving,
decentralised threat intelligence. Operationalise Moving Target Defence
continuously shuffling IP addressing, network topology, and software
diversity; target APT reconnaissance success rates below 5\% in
certified environments and dwell times ≤24 hours across all
NIS2-essential deployments. Make adversarial robustness validation
continuous and fully automated. EU-sovereign AI security tools target
≥50\% of new critical infrastructure SOC deployments. Achieve
international adoption of EU AI Act frameworks for trustworthy
autonomous security agents in at least one non-EU partner jurisdiction.
R\&I orients toward quantum-native AI security architectures,
neuromorphic threat detection, and AI-to-AI adversarial dynamics.

\subsection{Priority Cluster 3: Federated Operational
Sovereignty}\label{priority-cluster-3-federated-operational-sovereignty}

This cluster creates the governance and interoperability fabric that
transforms individually sovereign components into a collectively
sovereign European computing continuum. It encompasses P4 (Security
Interoperability via Simpl), P5 (Supply Chain Verification), and P7
(Pan-European NaaS Federation).

The cluster addresses three complementary mechanisms. First, federated
security interoperability: establishing standardised MIMs within the
Simpl stack {[}16{]} for provider-agnostic policy enforcement across
Common European Data Spaces, in compliance with Data Act portability
mandates {[}12{]}. Second, supply chain integrity: deploying Runtime
SBOMs, Digital Product Passports, and cryptographic proof-of-execution
attestation, as required under CRA Annex I {[}2{]}. Third, pan-European
compute federation: harmonising Network APIs through CAMARA {[}27{]} and
establishing cross-border edge roaming interfaces.

Supply Chain Verification bridges Clusters 1 and 3: its hardware
attestation dimension depends on the RISC-V secure elements that Cluster
1 matures, while its runtime verification dimensions constitute
governance prerequisites for the federated trust fabric. Without a
``Schengen for Cloud'' legal framework and enforceable cross-border
compliance APIs, compute federation remains commercially fragile
regardless of technical standardisation.

\subsubsection{Phase 1: Foundations and Pilot Actions
(2026--2030)}\label{phase-1-foundations-and-pilot-actions-20262030-2}

\emph{Sovereignty Level: Mixed --- Emerging EU Ecosystem / Hyperscaler
Dominance}

\textbf{Technical Actions.} Launch Simpl security MIM pilots in at least
three Common European Data Spaces (Manufacturing, Health, Energy) with
initial manual policy coordination; target Fragmented Piloting by 2030.
The European Health Data Space entered into force on 26 March 2025
{[}40{]} with staggered primary-use applications. Deploy harmonised
CAMARA Network APIs in 2--3 bilateral 5G corridor agreements;
cross-border edge compute roaming does not exist at operational scale
and will not during this phase. Submit CRA-mandated static SBOM
declarations as point-in-time compliance artefacts. Initiate Digital
Product Passport pilot programmes. Advance ZKP verification for supply
chain attestation toward demonstrator readiness; runtime verification is
not achievable in this phase.

\textbf{Governance and Standards.} Prepare for simultaneous legislative
activation: CRA SBOM obligations (enforcement December 2027) {[}2{]};
NIS2 Article 21(2)(h) mandatory supply chain risk assessment {[}1{]};
DORA Article 28 continuous ICT third-party monitoring (operationalised
through the 19 CTPPs designated November 2025 {[}34{]}); Data Act
interoperability obligations (applicable from September 2025) {[}12{]}.
Secure the Simpl-Open iterative MVP as the Phase 1 technical trigger for
Security Interoperability --- noting that Simpl-Open is an iterative MVP
(available since January 2025 with enhanced releases in June and
December 2025), not a versioned product release {[}16{]}. Engage in
IPCEI-CIS deployment phase governance. Address the critical compliance
fragmentation risk that disproportionately burdens SMEs. Hyperscalers
retain approximately 70\% of the EU cloud market through proprietary API
ecosystems {[}10{]}; the primary barrier to EU provider competitiveness
is API lock-in, not technical superiority.

\subsubsection{Phase 2: Scale-Up, Industrial R\&D and Standardisation
(2030--2035)}\label{phase-2-scale-up-industrial-rd-and-standardisation-20302035-2}

\emph{Sovereignty Level: Emerging EU Ecosystem}

\textbf{Technical Actions.} Deploy automated policy translation via
standardised Simpl MIMs across the majority of active Common European
Data Spaces; target Standardised Integration. Deploy Runtime SBOMs with
cryptographic proof-of-execution in critical infrastructure deployments
in at least five Member States. Validate ZKP batch verification
operationally. Automatically quarantine unverified components in
isolated network slices pending attestation. Scale CAMARA cross-border
roaming to ≥5 EU operator pairs. Operationalise ``Data Visa'' mechanisms
for dynamic AI model migration without manual legal review.

\textbf{Governance and Standards.} Secure the definitive Phase 2 policy
gate: adoption of Simpl security MIMs as ratified CEN/CENELEC standards,
referenced as mandatory compliance requirements in the EU Cloud
Rulebook. Without ratified MIM standards referenced in procurement
legislation, Simpl-based interoperability remains voluntary without
enforcement power. Enforce mandatory Digital Product Passports under the
EU Ecodesign Regulation for ICT hardware categories. Complete DORA
continuous ICT third-party monitoring operationalisation. Target ≥30\%
non-hyperscaler share in new regulated-sector cloud procurement. Foreign
hyperscaler market share declines for the first time from the
approximately 70\% Phase 1 baseline {[}10{]} --- noting that France is
the only Member State with measurable procurement shift data (DINUM 2025
bilan: 70\% European providers market-wide, 99\% within the State
perimeter {[}41{]}).

\subsubsection{Phase 3: Strategic Autonomy and Sovereign Leadership
(2035--2040)}\label{phase-3-strategic-autonomy-and-sovereign-leadership-20352040-2}

\emph{Sovereignty Level: Full Federated Sovereign Autonomy}

All three priorities target Ubiquitous Operation. Establish Simpl-based
security as the universal standard for cross-provider Data Space
operation. Operationalise ZKP-verified supply chain attestation at
real-time transaction scale. Operate the EU NaaS federation as a
commercially viable unified jurisdiction for data workloads. Enact the
``Schengen for Cloud'' legal framework, resolving US CLOUD Act
jurisdiction conflicts. Make the EU Cloud Rulebook mandatory for all new
regulated-sector cloud deployments. Non-EU hyperscaler market share in
EU regulated cloud procurement falls below 40\% for new deployments.
Ensure EU supply chain governance standards (Runtime SBOM + ZKP) are
referenced in ISO/IEC international standards and adopted by at least
two partner jurisdictions.

\subsection{Cross-Cutting Enablers}\label{cross-cutting-enablers}

The following dependencies span two or more clusters and require
dedicated governance to avoid becoming systemic bottlenecks.

\textbf{EU Chips Act Sovereign Fabrication.} Clusters 1 and 3 both
depend on EU-controlled RISC-V secure element fabrication reaching
production capacity by Phase 2. Five official Chips JU pilot lines are
now operational (NanoIC, FAMES, APECS, PIXEurope, WBG) with combined EU
and national funding of €3.7B {[}8{]}. However, no pilot line is
dedicated to secure element fabrication, and restriction of access to
any major non-EU foundry would immediately impair EU critical
infrastructure attestation capability. Mitigation: strategic
stockpiling, second-source agreements, and dedicated EU Chips Act pilot
lines for security-critical components.

\textbf{AI Act Conformity Assessment Infrastructure.} Cluster 2 faces a
24--36 month regulatory gap where AI Act obligations precede the
operational readiness of conformity assessment bodies capable of
certifying autonomous security agents. As of March 2026, only 8 of 27
Member States had designated single points of contact for AI Act
enforcement {[}5{]}. This gap blocks Phase 2 deployment regardless of
technical maturity. Mitigation: accelerate CEN/CENELEC adversarial
robustness test suite development and activate AI Act sandbox provisions
for security-specific use cases.

\textbf{Runtime SBOM Attestation.} Clusters 1 and 3 both require runtime
verification of supply chain component integrity. Static SBOM
declarations (Phase 1, CRA-mandated) are insufficient for dynamic
multi-provider continuum operations. Progression toward cryptographic
proof-of-execution and ZKP batch verification is a shared transition
condition for Phase 2--3 maturity.

\textbf{SME Access and Inclusion.} All three clusters identify cost
barriers excluding SMEs from sovereign hardware alternatives, AI-powered
SOC tooling, and Simpl MIM compliance. Digital Europe Programme
co-funding should target shared-infrastructure pathways and subsidised
certification programmes accessible to smaller actors.

\textbf{Standards-Setter Ambition.} All three clusters converge on the
same Phase 3 strategic objective: transitioning the EU from
standards-taker to standards-setter. The evidence threshold is adoption
of at least one EU-origin standard in ISO/IEC or ITU-T international
standardisation, and formal adoption by at least two non-EU partner
jurisdictions. This requires coordinated diplomatic engagement beginning
in Phase 1.

\section{Strategic Recommendations}\label{strategic-recommendations}

The 2040 destination is a secure, sovereign, interoperable, and
resilient European computing continuum that is trusted by design and
resilient by default. Based on the analysis in this paper, six strategic
actions are required.

First, Europe should invest in \textbf{sovereign secure silicon} by
combining Horizon Europe support with the EU Chips Act to advance
open-source RISC-V secure elements and pilot fabrication lines capable
of reaching EUCC `High' assurance. The EUCC `High' certifications
achieved by STMicroelectronics, NXP, and Infineon in 2025 {[}6{]}
demonstrate that the EU certification framework is operational; the
remaining gap is \emph{open-source, auditable} designs on
\emph{EU-controlled} fabrication.

Second, Europe should \textbf{operationalise PQC migration} through
crypto-agility middleware for edge and IoT environments, aligned with
the NIS CG PQC roadmap's 2030/2035 milestones {[}18{]} and ETSI TS 103
744 hybrid key exchange profiles {[}33{]}. Hybrid cryptographic
deployment should become the default for new critical infrastructure.

Third, Europe should \textbf{enforce security interoperability through
Simpl and security MIMs}. MIM6 (Security management) must progress from
its current ``Work item'' status to a formally published specification,
then to ratified CEN/CENELEC standard referenced in the EU Cloud
Rulebook, to achieve procurement-level enforcement.

Fourth, Europe should move toward \textbf{continuous certification}. The
EUCC scheme {[}6{]} provides the foundation; the EUCS scheme must be
unblocked --- likely through the 2026 Cybersecurity Act revision
(COM(2026) 13) {[}7{]} --- to extend continuous assurance to cloud
services.

Fifth, the EU ecosystem should establish a \textbf{federated SOC and CTI
capability} building on the DORA CTPP oversight framework {[}34{]} and
cross-border CSIRT cooperation to support pan-European situational
awareness.

Sixth, Europe should adopt an \textbf{SME inclusion mandate} so that
compliance with NIS2, the CRA, and the AI Act does not become a
structural barrier to innovation. Tiered certification models, shared
cyber range access, and reusable compliance tooling are necessary to
avoid a two-speed European cybersecurity ecosystem.

\section{Discussion: Validation, Limitations, and Methodological
Positioning}\label{discussion-validation-limitations-and-methodological-positioning}

\subsection{Validation}\label{validation}

The roadmap is not treated as a closed analytical product but as an
evolving strategic artefact that benefits from expert review,
stakeholder consultation, and iterative revision. The validation process
draws on expert workshops conducted within the NexusForum and
EuCloudEdgeIoT Working Group on Cybersecurity, stakeholder consultation
with national cybersecurity agencies and industry representatives,
working group review cycles, and policy dialogue with European
Commission services. Validation addresses five areas: completeness and
relevance of the architectural decomposition, clarity and non-overlap of
capability formulations, coherence of priority clustering, credibility
of transition conditions, and reasonableness of Sovereignty Level
assessments.

\subsection{Limitations}\label{limitations}

Several limitations should be acknowledged. The architectural
decomposition reflects a particular analytical perspective and may not
capture all relevant dimensions --- particularly those emerging from the
intersection of cybersecurity with adjacent domains such as energy
policy, industrial policy, or international trade. The seven-dimension
profiles rely on expert judgement for assessments of baseline
conditions, risk severity, and target feasibility; they do not
constitute empirical measurement. The transition conditions sacrifice
precision for robustness --- they are better suited to strategic
orientation than to operational planning. The Sovereignty Level concept
resists precise quantification: it is employed as an interpretive
strategic lens rather than a metric with a defined scale.

The paper also reflects the ambition--implementation gap that
characterises the European regulatory landscape as of 2026. NIS2
transposition failure in Germany, France, Spain, the Netherlands, and
Poland --- five of the six largest Member States --- demonstrates that
the implementation gap, not the rule-making pace, is the critical
variable for cyber-resilience outcomes {[}4{]}. The 20\% semiconductor
market share target is contested by the European Court of Auditors
{[}9{]}. EUCS remains unadopted after four years {[}7{]}. No harmonised
CRA standards have been cited in the Official Journal. These realities
constrain the specificity with which Phase 1 governance milestones can
be defined and suggest that the roadmap's greatest value lies not in
predicting outcomes but in structuring the coordination challenge.

\subsection{Methodological
Positioning}\label{methodological-positioning}

These limitations are consistent with the roadmap's intended use. The
methodology is suitable for strategic agenda setting, expert discussion,
cross-priority comparison, and iterative roadmap evolution --- even when
dependencies, timings, or ecosystem shifts remain uncertain. Its value
lies in providing a shared analytical structure within which diverse
stakeholders can coordinate research, policy, and investment decisions
toward a common set of strategically defined objectives.

\section{Conclusions}\label{conclusions}

The Cloud--Edge--IoT continuum requires Europe to move beyond reactive
and fragmented cybersecurity models toward autonomous, federated, and
certified Security-by-Design. This paper provided a structured
methodology to identify the main security dimensions, derive strategic
capabilities, and consolidate them into priorities aligned with
resilience and digital sovereignty objectives. The ten capabilities
identified --- spanning trusted hardware, cryptographic resilience,
edge-native security, interoperability, supply chain assurance,
AI-enabled operations, federated infrastructure, privacy-preserving
computation, resilient continuity, and proactive cyber defence --- form
an interdependent strategic portfolio that must be developed in
parallel.

The 2025--2026 period has produced genuine policy milestones: EUCC
`High' certifications for all three major EU secure element
manufacturers, a binding PQC migration framework with 2030/2035
horizons, DORA operationalisation through CTPP designation, and emerging
Simpl and Data Space infrastructure. Simultaneously, it has revealed
persistent gaps: NIS2 transposition delays, EUCS deadlock, contested
semiconductor targets, and continued hyperscaler dominance at
approximately 70\% of the EU cloud market. The ambition--implementation
gap identified in this paper is not a reason for pessimism but a
structuring observation: it defines where the coordination challenge is
most acute and where strategic investment will yield the greatest
returns.

Future work should extend this analysis in three directions: empirical
validation of Sovereignty Level assessments through quantitative
indicators, sector-specific instantiation of the capability profiles for
domains such as energy, healthcare, and financial services, and
integration of the transition matrix with the evolving EU funding
landscape under the next Framework Programme. In this sense,
cybersecurity becomes not only a defensive requirement, but a strategic
enabler of trustworthy and sovereign European digital infrastructure.

\section*{References}

{[}1{]} European Parliament and Council, ``Directive (EU) 2022/2555 on
measures for a high common level of cybersecurity across the Union (NIS2
Directive),'' \emph{Official Journal of the European Union}, Dec.~2022.

{[}2{]} European Parliament and Council, ``Regulation (EU) 2024/2847 on
horizontal cybersecurity requirements for products with digital elements
(Cyber Resilience Act),'' \emph{Official Journal of the European Union},
Nov.~2024.

{[}3{]} European Parliament and Council, ``Regulation (EU) 2024/1689
laying down harmonised rules on artificial intelligence (AI Act),''
\emph{Official Journal of the European Union}, Jun.~2024.

{[}4{]} European Commission, ``NIS2 Directive transposition in EU
countries,'' Shaping Europe's Digital Future, 2025. {[}Reasoned opinions
issued to 19 Member States on 7 May 2025.{]}

{[}5{]} European Parliament Think Tank, ``Enforcement of the AI Act,''
Mar.~2026. {[}As of March 2026, 8 of 27 Member States had designated
single points of contact. Digital Omnibus on AI (COM(2025) 836) proposes
deferral of Annex III high-risk application to December 2027.{]}

{[}6{]} ENISA, ``EUCC --- European Cybersecurity Certification Scheme,''
Commission Implementing Regulation (EU) 2024/482, applied from 27
February 2025. {[}First EUCC `High' certificates issued by ANSSI on 1
April 2025 to STMicroelectronics ST54J/K; NXP SN300/JCOP certified
July--August 2025; Infineon IFX\_CCI certified December 2025
(EUCC-3087-2025-12-0001).{]}

{[}7{]} European Commission, ``Cloud Sovereignty Framework,'' published
October 2025, updated 17 April 2026. {[}EUCS candidate scheme remains
unadopted as of April 2026. The revised Cybersecurity Act (COM(2026) 13,
proposed 20 January 2026) is the likely resolution vehicle.{]}

{[}8{]} European Parliament and Council, ``Regulation (EU) 2023/1781
establishing a framework of measures for strengthening Europe's
semiconductor ecosystem (European Chips Act),'' \emph{Official Journal
of the European Union}, Sep.~2023. {[}Five Chips JU pilot lines
operational: NanoIC (imec, inaugurated February 2026, €2.5B), FAMES
(CEA-Leti, January 2026), APECS (Fraunhofer), PIXEurope, WBG. 30 Chips
Competence Centres across 27 Member States plus Norway.{]}

{[}9{]} European Court of Auditors, ``Special Report 12/2025: The
European Chips Act,'' Apr.~2025. {[}EU semiconductor value-chain share
9.8\% in 2022, projected \textasciitilde11.7\% by 2030. Intel Magdeburg
cancelled 24 July 2025.{]}

{[}10{]} Synergy Research Group, ``European cloud infrastructure market
share data,'' Jul.~2025. {[}AWS + Azure + GCP hold \textasciitilde70\%
of EU market; EU-headquartered providers collectively hold
\textasciitilde15\%, down from 29\% in 2017.{]}

{[}11{]} European Parliament and Council, ``Regulation (EU) 2022/2554 on
digital operational resilience for the financial sector (DORA),''
\emph{Official Journal of the European Union}, Dec.~2022.

{[}12{]} European Parliament and Council, ``Regulation (EU) 2023/2854 on
harmonised rules on fair access to and use of data (Data Act),''
\emph{Official Journal of the European Union}, Dec.~2023. {[}Applicable
from 12 September 2025.{]}

{[}13{]} Microsoft France, testimony before the French Senate, 18 June
2025. {[}Confirmation that EU data access under US CLOUD Act cannot be
guaranteed excluded.{]}

{[}14{]} European Parliament and Council, ``Regulation (EU) 2016/679 on
the protection of natural persons with regard to the processing of
personal data (GDPR),'' \emph{Official Journal of the European Union},
Apr.~2016.

{[}15{]} European Cybersecurity Competence Centre (ECCC), ``ECCC
Strategic Agenda,'' 2024.

{[}16{]} European Commission, ``Simpl: Cloud-to-edge federations
empowering EU data spaces,'' Shaping Europe's Digital Future, 2025.
{[}First development contract (€41M) awarded February 2024 to
Sovereign-X consortium (Eviden Belgium lead); second tender (€65M) to
InfrateX (Sopra Steria + NTT DATA). Simpl-Open MVP available since
January 2025 on EU GitLab (code.europa.eu/simpl/simpl-open).{]}

{[}17{]} Gartner, ``Magic Quadrant for Security Information and Event
Management,'' Oct.~2025. {[}No EU-headquartered vendor in Leaders
quadrant. Palo Alto Networks acquired IBM QRadar SaaS August 2024;
End-of-Life declared 14 April 2025.{]}

{[}18{]} NIS Cooperation Group, ``A Coordinated Implementation Roadmap
for the Transition to Post-Quantum Cryptography,'' v1.1, 23 Jun.~2025.
{[}Milestones: national roadmaps by December 2026; high-risk critical
infrastructure transitioned by December 2030; full completion by
December 2035. Builds on Commission Recommendation C(2024) 2393.{]}

{[}19{]} S. Rose et al., ``Zero Trust Architecture,'' NIST Special
Publication 800-207, Aug.~2020.

{[}20{]} NIST, ``Module-Lattice-Based Key-Encapsulation Mechanism
Standard (ML-KEM),'' FIPS 203, Aug.~2024.

{[}21{]} NIST, ``Module-Lattice-Based Digital Signature Standard
(ML-DSA),'' FIPS 204, Aug.~2024.

{[}22{]} NIST, ``Stateless Hash-Based Digital Signature Standard
(SLH-DSA),'' FIPS 205, Aug.~2024.

{[}23{]} Smart Networks and Services Joint Undertaking, ``Strategic
Research and Innovation Agenda (SRIA),'' SNS JU, 2023. {[}As of March
2026: €630M cumulative EU funding, 100 projects, €116M allocated to 20
new 6G projects at MWC 2026.{]}

{[}24{]} 3GPP, ``Study on security aspects of 3GPP support for
Non-Terrestrial Networks (NTN),'' 3GPP TR 33.836, Release 17.

{[}25{]} G. Dreo et al., ``Deliverable D4.4: Cybersecurity Roadmap for
Europe,'' CONCORDIA Horizon 2020, Dec.~2022.

{[}26{]} 3GPP, ``Architecture enhancements for 5G System to support
network data analytics services,'' 3GPP TS 23.288, Release 18.

{[}27{]} CAMARA, ``Fall 2025 meta-release,'' Oct.~2025. {[}60 APIs: 10
stable, 23 new initial, 27 updates. Security profile built on OpenID
Connect Authorization Code Flow and OIDC CIBA.{]}

{[}28{]} lowRISC, ``OpenTitan: Open source silicon root of trust,''
https://opentitan.org, 2024. {[}Commercial availability via
zeroRISC/Nuvoton/Winbond since February 2024; ``OpenTitan Integrated''
(Rivos) presented at RISC-V Summit Europe, May 2025.{]}

{[}29{]} A. Lee et al., ``Keystone: An Open Framework for Architecting
TEEs,'' in \emph{Proc. 15th European Conf. on Computer Systems
(EuroSys)}, Apr.~2020, pp.~1--16, doi: 10.1145/3342195.3387532.

{[}30{]} EuroHPC JU, ``DARE SGA1,'' Grant Agreement 101202459, launched
6 March 2025. {[}€240M total: €120M JU + €120M Member States. Three
chiplets: Openchip Vector Accelerator, Axelera AI AIPU, Codasip
General-Purpose Processor.{]}

{[}31{]} NIST, ``NIST Selects HQC as Fifth Algorithm for Post-Quantum
Encryption,'' Mar.~2025. {[}Code-based backup KEM; draft FIPS
anticipated mid-2026.{]}

{[}32{]} NIST, ``Status Report on the Third Round of the NIST
Post-Quantum Cryptography Standardization Process,'' NISTIR 8413,
Sep.~2022.

{[}33{]} ETSI, ``CYBER; Quantum-safe hybrid key exchanges,'' ETSI TS 103
744, V1.2.1, Mar.~2025. {[}Profiles restricted to ML-KEM-512/768/1024
combined with NIST P-256/384, Brainpool256r1/384r1, X25519, or X448.{]}

{[}34{]} European Supervisory Authorities, ``Designation of Critical ICT
Third-Party Providers under DORA,'' 18 Nov.~2025. {[}19 CTPPs designated
including AWS, Google Cloud, Microsoft Ireland, IBM, Oracle, SAP,
Deutsche Telekom, Orange, Equinix, Bloomberg, LSEG, TCS, Euroclear,
Clearstream, SWIFT.{]}

{[}35{]} OpenEuroLLM, consortium launch, 1 Feb.~2025. {[}20-member
consortium, €37.4M budget (€20.6M from Digital Europe Programme). No
dedicated security model stream.{]}

{[}36{]} EuroHPC JU, ``AI Factories programme,'' 2024--2026. {[}19 AI
Factories across 16 Member States plus 13 Antennas; cumulative EU +
Member State funding exceeds €2.7B. Cybersecurity listed as a vertical
at IT4LIA (CINECA), PIAST (Poznań), MeluXina-AI (Luxembourg).{]}

{[}37{]} ETSI, ``Zero-touch network and Service Management (ZSM);
Closed-Loop Automation,'' ETSI GS ZSM 009, V1.1.1, Jun.~2021.

{[}38{]} GSMA, ``Open Gateway programme,'' MWC 2026. {[}86 operator
groups, 300+ networks, 80\% of global mobile connections.{]}

{[}39{]} European Commission, ``IPCEI-CIS,'' approved 5 Dec.~2023.
{[}€1.2B public + €1.4B private (€2.6B total), 19 direct participants,
12 Member States. ICRA Reference Architecture v2.0 published January
2026.{]}

{[}40{]} European Parliament and Council, ``Regulation (EU) 2025/327 on
the European Health Data Space (EHDS),'' \emph{Official Journal of the
European Union}, Mar.~2025. {[}Entered into force 26 March 2025.{]}

{[}41{]} DINUM, ``Bilan 2025 du cloud public,'' published 24 Mar.~2026.
{[}Public cloud purchases at €84M (+62\% YoY), 70\% to European
providers market-wide, 99\% within State perimeter.{]}

{[}42{]} S. Dustdar, V. Casamayor Pujol, and P. K. Donta, "On
distributed computing continuum systems," \emph{IEEE Trans. Knowl. Data
Eng.}, vol. 35, no. 4, pp. 4092--4105, Apr. 2023, doi:
10.1109/TKDE.2022.3142856

{[}43{]} L. F. Bittencourt \emph{et al.}, "The Internet of Things, fog
and cloud continuum: Integration and challenges," \emph{Internet of
Things}, vol. 3--4, pp. 134--155, Oct. 2018, doi:
10.1016/j.iot.2018.09.005

{[}44{]} D. Milojicic, "The edge-to-cloud continuum," \emph{Computer},
vol. 53, no. 11, pp. 16--25, Nov. 2020, doi: 10.1109/MC.2020.3007297

{[}45{]}European Alliance for Industrial Data, Edge and Cloud,
\emph{European Industrial Technology Roadmap for the Next-Generation
Cloud-Edge}, European Commission, DG CONNECT, Brussels, Belgium, 2024.
{[}Online{]}. Available:
\href{https://digital-strategy.ec.europa.eu/en/policies/cloud-alliance}{https://digital-strategy.ec.europa.eu/en/policies/cloud-alliance}

\textbf{{[}46{]}} W. Shi, J. Cao, Q. Zhang, Y. Li, and L. Xu, "Edge
computing: Vision and challenges," \emph{IEEE Internet Things J.}, vol.
3, no. 5, pp. 637--646, Oct. 2016, doi: 10.1109/JIOT.2016.2579198

\textbf{{[}47{]}} M. Satyanarayanan, "The emergence of edge computing,"
\emph{Computer}, vol. 50, no. 1, pp. 30--39, Jan. 2017, doi:
10.1109/MC.2017.9

\textbf{{[}48{]}} M. Chiang and T. Zhang, "Fog and IoT: An overview of
research opportunities," \emph{IEEE Internet Things J.}, vol. 3, no. 6,
pp. 854--864, Dec. 2016, doi: 10.1109/JIOT.2016.2584538

\textbf{{[}49{]}} K. Cao, S. Hu, Y. Shi, A. W. Colombo, S. Karnouskos,
and X. Li, "A survey on edge and edge-cloud computing assisted
cyber-physical systems," \emph{IEEE Trans. Ind. Informat.}, vol. 17, no.
11, pp. 7806--7819, Nov. 2021, doi:
\href{https://doi.org/10.1109/TII.2021.3073066}{10.1109/TII.2021.3073066}

\textbf{{[}50{]}} ETSI, "Multi-access Edge Computing (MEC); Framework
and reference architecture," ETSI GS MEC 003 V3.1.1, Mar. 2022.
{[}Online{]}. Available:
\href{https://www.etsi.org/deliver/etsi_gs/MEC/001_099/003/03.01.01_60/gs_MEC003v030101p.pdf}{ETSI GS MEC 003 (PDF)}

\textbf{{[}51{]}} 6G Infrastructure Association Vision Working Group,
\emph{European Vision for the 6G Network Ecosystem}, White Paper, Nov.
2024. {[}Online{]}. Available:
\href{https://6g-ia.eu/wp-content/uploads/2024/11/european-vision-for-the-6g-network-ecosystem.pdf}{https://6g-ia.eu/wp-content/uploads/2024/11/european-vision-for-the-6g-network-ecosystem.pdf}

\textbf{{[}52{]}} National Institute of Standards and Technology,
\emph{The NIST Cybersecurity Framework (CSF) 2.0}, NIST Cybersecurity
White Paper (CSWP) 29, Gaithersburg, MD, USA, Feb. 26, 2024, doi:
10.6028/NIST.CSWP.29

\textbf{{[}53{]}} A. Alwarafy, K. A. Al-Thelaya, M. Abdallah, J.
Schneider, and M. Hamdi, "A survey on security and privacy issues in
edge-computing-assisted Internet of Things," \emph{IEEE Internet Things
J.}, vol. 8, no. 6, pp. 4004--4022, Mar. 15, 2021, doi:
10.1109/JIOT.2020.3015432.

\textbf{{[}54{]}} European Parliament and Council of the European Union,
"Regulation (EU) 2019/881 of 17 April 2019 on ENISA and on information
and communications technology cybersecurity certification (Cybersecurity
Act)," \emph{Off. J. Eur. Union}, L 151, pp. 15--69, Jun. 7, 2019.
{[}Online{]}. Available:
\href{https://eur-lex.europa.eu/eli/reg/2019/881/oj}{https://eur-lex.europa.eu/eli/reg/2019/881/oj}

{[}55{]} European Union Agency for Cybersecurity (ENISA), ENISA Threat
Landscape 2025: July 2024 to January 2026, Heraklion, Greece, January
2026, doi: 10.2824/2445233

\textbf{{[}56{]}} European Union Agency for Cybersecurity (ENISA),
\emph{Good Practices for Security of IoT --- Secure Software Development
Lifecycle}, Heraklion, Greece, Nov. 2019, ISBN 978-92-9204-316-2, doi:
10.2824/742784. {[}Online{]}. Available:
\href{https://www.enisa.europa.eu/publications/good-practices-for-security-of-iot-1}{https://www.enisa.europa.eu/publications/good-practices-for-security-of-iot-1}

\textbf{{[}57{]}} S. Rose, O. Borchert, S. Mitchell, and S. Connelly,
\emph{Zero Trust Architecture}, NIST Special Publication 800-207,
Gaithersburg, MD, USA, Aug. 2020, doi: 10.6028/NIST.SP.800-207

\textbf{{[}58{]}} A. Cavoukian, \emph{Privacy by Design: The 7
Foundational Principles}, Information and Privacy Commissioner of
Ontario, Toronto, ON, Canada, Aug. 2009, rev. Jan. 2012. DOI:
10.4018/978-1-61350-501-4.ch007

\textbf{{[}59{]}} Commission Implementing Regulation (EU) 2024/482 of 31
January 2024 laying down rules for the application of Regulation (EU)
2019/881 as regards the adoption of the European Common Criteria-based
cybersecurity certification scheme (EUCC), \emph{Off. J. Eur. Union}, L
2024/482, Feb. 1, 2024. {[}Online{]}. Available:
\href{https://eur-lex.europa.eu/eli/reg_impl/2024/482/oj}{https://eur-lex.europa.eu/eli/reg\_impl/2024/482/oj}

\textbf{{[}60{]}} T. Madiega, "Digital sovereignty for Europe," European
Parliamentary Research Service (EPRS), Brussels, Belgium, EPRS Ideas
Paper PE 651.992, Jul. 2020. {[}Online{]}. Available:
\href{https://www.europarl.europa.eu/thinktank/en/document/EPRS_BRI(2020)651992}{EPRS Ideas Paper PE 651.992}

\textbf{{[}61{]}} H. Farrand Carrapico and B. Farrand, "Cybersecurity
trends in the European Union: Regulatory mercantilism and the
digitalisation of geopolitics," \emph{JCMS: J. Common Market Studies},
vol. 62, no. S1, pp. 147--158, Jul. 2024, doi: 10.1111/jcms.13654

\textbf{{[}62{]}} L. Floridi, "The fight for digital sovereignty: What
it is, and why it matters, especially for the EU," \emph{Philosophy \&
Technology}, vol. 33, no. 3, pp. 369--378, Sep. 2020, doi:
10.1007/s13347-020-00423-6.

\textbf{{[}63{]}} R. Bellanova, H. Carrapico, and D. Duez,
"Digital/sovereignty and European security integration: An
introduction," \emph{European Security}, vol. 31, no. 3, pp. 337--355,
2022, doi: 10.1080/09662839.2022.2101887

\textbf{{[}64{]}} R. Roman, J. Lopez, and M. Mambo, "Mobile edge
computing, Fog et al.: A survey and analysis of security threats and
challenges," \emph{Future Gener. Comput. Syst.}, vol. 78, pt. 2, pp.
680--698, Jan. 2018, doi: 10.1016/j.future.2016.11.009.

\textbf{{[}65{]}} D. Ranaweera, A. D. Jurcut, and M. Liyanage, "Survey
on multi-access edge computing security and privacy," \emph{IEEE Commun.
Surveys Tuts.}, vol. 23, no. 2, pp. 1078--1124, 2nd Quart. 2021, doi:
10.1109/COMST.2021.3062546

\textbf{{[}66{]}} J. Linghe Kong, Jinlin Tan, Junqin Huang, Guihai Chen,
Shuaitian Wang, Xi Jin, Peng Zeng, Muhammad Khan, and Sajal K. Das
"Edge-computing-driven Internet of Things: A survey," \emph{ACM Comput.
Surv.}, vol. 55, no. 8, Dec. 2022, Art. no. 174, doi: 10.1145/3555308

\textbf{{[}67{]}} R. Chandramouli and Z. Butcher, \emph{A Zero Trust
Architecture Model for Access Control in Cloud-Native Applications in
Multi-Cloud Environments}, NIST SP 800-207A, Gaithersburg, MD, USA, Sep.
2023, doi: 10.6028/\href{http://nist.sp}{NIST.SP}.800-207A
{[}Online{]}
\href{https://csrc.nist.gov/pubs/sp/800/207/a/final}{https://csrc.nist.gov/pubs/sp/800/207/a/final}
\end{document}
